\section{Τεχνική Ανάλυση του \ENG{SCORM} 2004}
\label{sec:scorm2004_analysis}
\label{sec:scorm2004_technical}

Η μετάβαση από το \ENG{SCORM} 1.2 στο \ENG{SCORM} 2004 συνιστά ουσιαστική αναδιαμόρφωση της τεχνικής φιλοσοφίας του προτύπου και όχι απλή αναθεώρηση ονοματολογίας. Το νέο πρότυπο επιδιώκει να διατηρήσει τη βασική λογική διαλειτουργικού πακεταρίσματος και ανταλλαγής δεδομένων, αλλά ταυτόχρονα εισάγει επίσημη προδιαγραφή \ENG{sequencing} και \ENG{navigation}, πλουσιότερο μοντέλο δεδομένων και αυστηρότερη σύνδεση ανάμεσα στη δομή περιεχομένου και στη συμπεριφορά εκτέλεσης~\cite{ADL_SCORM2004_CAM, ADL_SCORM_Users_Guide, IMS_SimpleSequencing}.

Η παρούσα ενότητα εξετάζει το \ENG{SCORM} 2004 μέσα από τα τρία βασικά του συστατικά, δηλαδή το \ENG{Content Aggregation Model} (\ENG{CAM}), το \ENG{Run-Time Environment} (\ENG{RTE}) και την προδιαγραφή \ENG{Sequencing and Navigation} (\ENG{SN}). Η σύγκριση με το \ENG{SCORM} 1.2 είναι αναπόφευκτη, επειδή πολλές από τις σχεδιαστικές επιλογές του \ENG{SCORM} 2004 γίνονται κατανοητές μόνο ως απάντηση στους περιορισμούς του προγενέστερου προτύπου. Ο Πίνακας~\ref{tab:scorm2004_vs_12} συνοψίζει τα κύρια σημεία διαφοροποίησης~\cite{SCORM_Comparison_Guide, Rustici_SCORM_Versions}.

\begin{table}[htbp]
\centering
\caption{Σύγκριση \ENG{SCORM} 1.2 και \ENG{SCORM} 2004}
\label{tab:scorm2004_vs_12}
\begingroup
\small
\setlength{\tabcolsep}{3pt}
\renewcommand{\arraystretch}{1.16}
\begin{adjustbox}{center,max width=\textwidth}
\begin{tabular}{|p{4cm}|p{5cm}|p{5cm}|}
\hline
Χαρακτηριστικό & \ENG{SCORM} 1.2 & \ENG{SCORM} 2004 \\ 
\hline
\hline
Υπογραφή \ENG{API} & \texttt{\ENG{LMSInitialize}()}, \texttt{\ENG{LMSFinish}()}, \texttt{\ENG{LMSGetValue}()}, \texttt{\ENG{LMSSetValue}()} & \texttt{\ENG{Initialize}()}, \texttt{\ENG{Terminate}()}, \texttt{\ENG{GetValue}()}, \texttt{\ENG{SetValue}()} \\ 
\hline
\ENG{Sequencing} & Περιορισμένο; μόνο επιλογή χρήστη & Προηγμένο \ENG{IMS SS}-βασισμένο \ENG{sequencing} με κανόνες \\ 
\hline
\ENG{Navigation} & Βασικό; χωρίς επίσημο μοντέλο & Πλήρες μοντέλο \ENG{navigation} με \ENG{events} \\ 
\hline
Μοντέλο δεδομένων & \texttt{\ENG{cmi.core}.*}, περιορισμένες \ENG{interactions} & Εκτεταμένη ιεραρχία \texttt{\ENG{cmi}.*}, πλουσιότερες \ENG{interactions} \\ 
\hline
\ENG{Suspend} \ENG{Data} & 4096 χαρακτήρες & 64000 χαρακτήρες (3\ENG{rd/4th} \ENG{Ed}.) \\ 
\hline
\ENG{Rollup} & Καθορίζεται από το \ENG{LMS} & Καθορίζεται από το περιεχόμενο μέσω \ENG{rollup} κανόνων \\ 
\hline
\ENG{Objectives} & Βασική παρακολούθηση & Καθολικά \ENG{objectives} με \ENG{mapping} \\ 
\hline
\ENG{Progress} \ENG{Measure} & Δεν υποστηρίζεται & \texttt{\ENG{cmi.progress}\_\ENG{measure}} \\ 
\hline
\ENG{Completion} \ENG{Threshold} & Δεν υποστηρίζεται & \texttt{\ENG{cmi.completion}\_\ENG{threshold}} \\ 
\hline
\end{tabular}
\end{adjustbox}
\endgroup

\end{table}

Η αρχιτεκτονική επισκόπηση του προτύπου, όπως παρουσιάζεται στο Σχήμα~\ref{fig:scorm2004_architecture}, αναδεικνύει τη μετάβαση από ένα σχετικά επίπεδο μοντέλο εκτέλεσης σε μια πολυεπίπεδη αρχιτεκτονική όπου το πακέτο περιεχομένου, το μοντέλο δραστηριοτήτων, το \ENG{runtime} και ο μηχανισμός \ENG{sequencing} συνεργάζονται στενά. Η προσθήκη \ENG{activity trees}, ελέγχων \ENG{navigation} και διαδικασιών \ENG{rollup} καθιστά το \ENG{SCORM} 2004 περισσότερο εκφραστικό, αλλά και σαφώς πιο απαιτητικό ως προς την υλοποίηση~\cite{ADL_SCORM2004_Testing, SCORM_Technical_Overview}.

\begin{figure}[htbp]
\centering
% Figure 3.6: SCORM 2004 Architectural Overview
\begin{chapterfigurebox}
\begin{tikzpicture}[
    node distance=1.1cm and 1.2cm,
    component/.style={rectangle, draw=blue!60, fill=blue!10, thick, minimum width=3cm, minimum height=1cm, align=center, rounded corners},
    subcomponent/.style={rectangle, draw=green!60, fill=green!10, thick, minimum width=2.5cm, minimum height=0.8cm, align=center},
    dataflow/.style={->, >=stealth, thick, blue!70},
    interaction/.style={<->, >=stealth, thick, red!70, dashed},
    layer/.style={rectangle, draw=gray!50, fill=gray!5, thick, minimum width=11.8cm, minimum height=2.3cm, align=center, rounded corners}
]

% Layer 1: LMS Infrastructure
\node[layer] (lms_layer) at (0,6) {\ENG{LMS} \ENG{Infrastructure} \ENG{Layer}};
\node[component, below=0.3cm of lms_layer] (lms_core) {\ENG{LMS} \ENG{Core}};
\node[component, left=of lms_core] (user_mgmt) {\ENG{User}\\\ENG{Management}};
\node[component, right=of lms_core] (reporting) {\ENG{Reporting} \&\\\ENG{Analytics}};

% Layer 2: SCORM 2004 Runtime Components
\node[layer, below=3cm of lms_layer] (runtime_layer) {\ENG{SCORM} 2004 \ENG{Runtime} \ENG{Layer}};
\node[component, below=0.3cm of runtime_layer] (rte_core) {\ENG{Run-Time}\\\ENG{Environment}\\(\ENG{RTE})};
\node[component, left=of rte_core] (api_adapter) {\ENG{API} \ENG{Adapter}\\(\ENG{API}\_1484\_11)};
\node[component, right=of rte_core] (sequencer) {\ENG{Sequencing}\\\ENG{Engine}};

% Layer 3: Content Components
\node[layer, below=3cm of runtime_layer] (content_layer) {\ENG{Content} \ENG{Layer}};
\node[component, below=0.3cm of content_layer] (content_pkg) {\ENG{Content}\\\ENG{Package}\\(\ENG{CAM})};
\node[subcomponent, left=of content_pkg, yshift=-0.5cm] (manifest) {\ENG{imsmanifest.xml}\\+ \ENG{Sequencing}};
\node[subcomponent, right=of content_pkg, yshift=-0.5cm] (sco) {\ENG{SCOs}\\+ \ENG{Assets}};

% Data flows
\draw[dataflow] (content_pkg) -- node[right, font=\scriptsize] {\ENG{Launch} \ENG{SCO}} (rte_core);
\draw[interaction] (sco) -- node[right, font=\scriptsize] {\ENG{API} \ENG{Calls}} (api_adapter);
\draw[dataflow] (api_adapter) -- (rte_core);
\draw[dataflow] (rte_core) -- node[right, font=\scriptsize] {\ENG{Tracking} \ENG{Data}} (lms_core);
\draw[dataflow] (manifest) -- node[left, font=\scriptsize, align=left] {\ENG{Sequencing}\\\ENG{Rules}} (sequencer);
\draw[dataflow] (sequencer) -- node[left, font=\scriptsize, align=left] {\ENG{Navigation}\\\ENG{Decisions}} (rte_core);
\draw[dataflow] (lms_core) -- (reporting);
\draw[dataflow] (user_mgmt) -- (lms_core);

% Activity Tree representation
\node[component, below=1.5cm of content_layer, xshift=-4cm, fill=yellow!10, draw=orange!60] (activity_tree) {\ENG{Activity} \ENG{Tree}\\(\ENG{Organizations})};
\draw[dataflow] (manifest) -- (activity_tree);
\draw[dataflow] (activity_tree) -- (sequencer);

% Annotations
\node[font=\scriptsize\itshape, text width=3cm, align=left, below=0.2cm of sco] {\ENG{Launch} \ENG{in}\\\ENG{browser} \ENG{window}};
\node[font=\scriptsize\itshape, text width=3cm, align=right, below=0.2cm of manifest] {\ENG{XML} \ENG{with}\\\ENG{IMS} \ENG{SS} \ENG{extensions}};

\end{tikzpicture}
\end{chapterfigurebox}

\caption{Αρχιτεκτονική Επισκόπηση \ENG{SCORM} 2004}
\label{fig:scorm2004_architecture}
\end{figure}

\subsection{Μοντέλο Συγκέντρωσης Περιεχομένου (\ENG{Content Aggregation Model} -- \ENG{CAM})}
\label{subsec:scorm2004_cam}

Το \ENG{SCORM} 2004 \ENG{CAM} κληρονομεί τον βασικό μηχανισμό πακεταρίσματος του \ENG{SCORM} 1.2, αλλά τον επεκτείνει ώστε το \ENG{manifest} να μετατραπεί από απλό δοχείο δομικής πληροφορίας σε φορέα δηλωτικής παιδαγωγικής λογικής. Η ενσωμάτωση στοιχείων από το \ENG{IMS Content Packaging} και το \ENG{IMS Simple Sequencing} σημαίνει ότι η περιγραφή του πακέτου περιλαμβάνει πλέον όχι μόνο τους πόρους και τη διάταξή τους, αλλά και κανόνες που επηρεάζουν ενεργά την πλοήγηση, την ολοκλήρωση και τη ροή μάθησης~\cite{ADL_SCORM2004_CAM, IMS_CP_Specification, ADL_SCORM2004_SN, SCORM_CAM_Extensions}.

\subsubsection{Ενισχυμένη δομή \ENG{manifest}}

Το \texttt{\ENG{imsmanifest.xml}} του \ENG{SCORM} 2004 διατηρεί τη βασική σύνταξη του \ENG{SCORM} 1.2, αλλά αποκτά σημαντικά μεγαλύτερη εκφραστικότητα μέσω πρόσθετων \ENG{namespaces} και στοιχείων \ENG{sequencing}. Το αποτέλεσμα είναι ένα \ENG{manifest} που δεν αρκείται στην περιγραφή της δομής περιεχομένου, αλλά ενσωματώνει δηλωτική συμπεριφορά η οποία επηρεάζει την παράδοση μαθησιακών δραστηριοτήτων. Το ενισχυμένο αυτό σχήμα παρουσιάζεται στο Σχήμα~\ref{fig:cam_2004_extensions}~\cite{ADL_SCORM2004_CAM, IMS_SS_InfoModel}.

\begin{figure}[htbp]
\centering
% Figure 3.7: SCORM 2004 CAM Manifest Extensions
\begin{chapterfigurebox}
\begin{tikzpicture}[
    every node/.style={font=\small},
    arrow/.style={->, >=latex, thick},
    required/.style={rectangle, draw=red!60, fill=red!10, thick, minimum width=3cm, minimum height=0.7cm, align=center, rounded corners},
    optional/.style={rectangle, draw=green!60, fill=green!10, thick, minimum width=3cm, minimum height=0.7cm, align=center, rounded corners}
]

\node[required] (manifest) at (0,0) {\texttt{<\ENG{manifest}>}\\\ENG{Root} \ENG{element}};

\node[optional] (metadata) at (-4.4,-2.2) {\texttt{<\ENG{metadata}>}\\\ENG{LOM} \ENG{metadata}};
\node[required] (organizations) at (0,-2.2) {\texttt{<\ENG{organizations}>}\\\ENG{Activity} \ENG{trees}};
\node[required] (resources) at (4.4,-2.2) {\texttt{<\ENG{resources}>}\\\ENG{Resource} \ENG{defs}};

\node[optional] (submanifest) at (-4.4,-4.4) {\texttt{<\ENG{submanifest}>}\\\ENG{Nested} \ENG{package}};
\node[required] (org) at (0,-4.4) {\texttt{<\ENG{organization}>}\\\ENG{Org} \ENG{structure}};
\node[required] (resource) at (4.4,-4.4) {\texttt{<\ENG{resource}>}\\\ENG{SCO/Asset}};

\node[required] (item) at (0,-6.5) {\texttt{<\ENG{item}>}\\\ENG{Activity}};
\node[optional] (file) at (2.8,-6.5) {\texttt{<\ENG{file}>}\\\ENG{Launch} \ENG{file}};
\node[optional] (dependency) at (6.0,-6.5) {\texttt{<\ENG{dependency}>}\\\ENG{Shared} \ENG{assets}};

\node[optional] (sequencing) at (-1.8,-8.6) {\texttt{<\ENG{imsss:sequencing}>}\\\ENG{Sequencing} \ENG{rules}};
\node[optional] (presentation) at (1.8,-8.6) {\texttt{<\ENG{adlnav:presentation}>}\\\ENG{Navigation} \ENG{controls}};

\draw[arrow] (manifest) -- (metadata);
\draw[arrow] (manifest) -- (organizations);
\draw[arrow] (manifest) -- (resources);
\draw[arrow] (manifest) -- (submanifest);
\draw[arrow] (organizations) -- (org);
\draw[arrow] (org) -- (item);
\draw[arrow] (item) -- (sequencing);
\draw[arrow] (item) -- (presentation);
\draw[arrow] (resources) -- (resource);
\draw[arrow] (resource) -- (file);
\draw[arrow] (resource) -- (dependency);

\node[required, minimum width=2.2cm] (leg_req) at (-1.4,-10.8) {\ENG{Required}};
\node[optional, minimum width=2.2cm] (leg_opt) at (1.4,-10.8) {\ENG{Optional}};
\node[font=\scriptsize\itshape, text width=6.3cm, align=center] at (0,-12.2) {\ENG{SCORM} 2004 \ENG{adds} \ENG{embedded} \ENG{sequencing} \ENG{rules} \ENG{and} \ENG{ADL} \ENG{navigation} \ENG{controls} \ENG{directly} \ENG{inside} \ENG{the} \ENG{manifest}.};

\end{tikzpicture}
\end{chapterfigurebox}

\caption{Επεκτάσεις \ENG{Manifest} στο \ENG{SCORM} 2004 \ENG{CAM}}
\label{fig:cam_2004_extensions}
\end{figure}

Σε δομικό επίπεδο, το \ENG{manifest} συνεχίζει να οργανώνεται γύρω από τα \ENG{metadata}, τα \ENG{organizations} και τα \ENG{resources}, αλλά σε καθένα από αυτά τα τμήματα προστίθενται νέες σημασιολογικές στρώσεις. Τα \ENG{metadata} υποστηρίζουν πληρέστερα προφίλ \ENG{LOM}, τα \ENG{organizations} μετασχηματίζονται σε ιεραρχικά \ENG{activity trees}, οι \ENG{resources} εξακολουθούν να δηλώνουν \ENG{assets} και \ENG{SCOs}, ενώ τα στοιχεία \ENG{sequencing} και \ENG{navigation} εισάγουν κανόνες που καθορίζουν τη ροή της μαθησιακής εμπειρίας. Ο Πίνακας~\ref{tab:cam_2004_elements} συνοψίζει τα κυριότερα συστατικά αυτού του ενισχυμένου \ENG{manifest}.

\begin{table}[htbp]
\centering
\caption{Στοιχεία \ENG{Manifest} του \ENG{SCORM} 2004 \ENG{CAM}}
\label{tab:cam_2004_elements}
\begingroup
\footnotesize
\setlength{\tabcolsep}{3pt}
\renewcommand{\arraystretch}{1.16}
\begin{adjustbox}{center,max width=\textwidth,max totalheight=0.86\textheight}
\begin{tabular}{|p{5.5cm}|p{8.5cm}|}
\hline
\ENG{Manifest} \ENG{Element} & \ENG{Description} \ENG{and} \ENG{Purpose} \\
\hline
\hline
\texttt{<\ENG{manifest}>} & \ENG{Root} \ENG{element}; \ENG{defines} \ENG{namespace} \ENG{declarations} (\texttt{\ENG{xmlns}}, \texttt{\ENG{xmlns:adlcp}}, \texttt{\ENG{xmlns:adlseq}}, \texttt{\ENG{xmlns:adlnav}}, \texttt{\ENG{xmlns:imsss}}) \\
\hline
\texttt{<\ENG{metadata}>} & \ENG{Package-level} \ENG{LOM} \ENG{metadata}; \ENG{includes} \ENG{schema} \ENG{name}, \ENG{schema} \ENG{version}, \ENG{and} \ENG{location} \\
\hline
\texttt{<\ENG{organizations}>} & \ENG{Container} \ENG{for} \ENG{one} \ENG{or} \ENG{more} \ENG{organization} \ENG{elements}; \ENG{defines} \ENG{activity} \ENG{trees} \\
\hline
\texttt{<\ENG{organization}>} & \ENG{Represents} \ENG{a} \ENG{complete} \ENG{learning} \ENG{sequence}; \ENG{corresponds} \ENG{to} \ENG{sequencing} \ENG{activity} \ENG{tree} \ENG{root} \\
\hline
\texttt{<\ENG{item}>} & \ENG{Activity} \ENG{within} \ENG{an} \ENG{organization}; \ENG{can} \ENG{be} \ENG{nested} \ENG{hierarchically}; \ENG{corresponds} \ENG{to} \ENG{sequencing} \ENG{activity} \\
\hline
\texttt{<\ENG{item} \ENG{identifier}>} & \ENG{Unique} \ENG{identifier} \ENG{for} \ENG{the} \ENG{item/activity} \ENG{within} \ENG{the} \ENG{manifest} \\
\hline
\texttt{<\ENG{item} \ENG{identifierref}>} & \ENG{References} \ENG{a} \texttt{<\ENG{resource}>} \ENG{element} \ENG{by} \ENG{identifier}; \ENG{empty} \ENG{for} \ENG{clusters} \\
\hline
\texttt{<\ENG{item} \ENG{isvisible}>} & \ENG{Boolean} \ENG{indicating} \ENG{whether} \ENG{item} \ENG{should} \ENG{appear} \ENG{in} \ENG{LMS} \ENG{TOC} \\
\hline
\texttt{<\ENG{imsss:sequencing}>} & \ENG{IMS} \ENG{Simple} \ENG{Sequencing} \ENG{rules} \ENG{embedded} \ENG{within} \texttt{<\ENG{item}>}; \ENG{defines} \ENG{sequencing} \ENG{behavior} \ENG{for} \ENG{the} \ENG{activity} \\
\hline
\texttt{<\ENG{imsss:controlMode}>} & \ENG{Sequencing} \ENG{control} \ENG{modes} (\ENG{choice}, \ENG{flow}, \ENG{forwardOnly}, \ENG{choiceExit}) \\
\hline
\texttt{<\ENG{imsss:sequencingRules}>} & \ENG{Container} \ENG{for} \ENG{precondition}, \ENG{postcondition}, \ENG{and} \ENG{exit} \ENG{rules} \\
\hline
\texttt{<\ENG{imsss:limitConditions}>} & \ENG{Attempt} \ENG{limits} \ENG{and} \ENG{duration} \ENG{limits} \ENG{for} \ENG{the} \ENG{activity} \\
\hline
\texttt{<\ENG{imsss:rollupRules}>} & \ENG{Defines} \ENG{rollup} \ENG{behavior} \ENG{for} \ENG{cluster} \ENG{activities} \\
\hline
\texttt{<\ENG{imsss:objectives}>} & \ENG{Local} \ENG{and} \ENG{global} \ENG{objectives} \ENG{associated} \ENG{with} \ENG{the} \ENG{activity} \\
\hline
\texttt{<\ENG{imsss:randomizationControls}>} & \ENG{Specifies} \ENG{randomization} \ENG{timing} \ENG{and} \ENG{whether} \ENG{to} \ENG{randomize} \ENG{children} \\
\hline
\texttt{<\ENG{imsss:deliveryControls}>} & \ENG{Tracked} \ENG{status}, \ENG{completion/objective} \ENG{set} \ENG{by} \ENG{content} \ENG{flags} \\
\hline
\texttt{<\ENG{adlnav:presentation}>} & \ENG{ADL} \ENG{navigation} \ENG{presentation} \ENG{settings}; \ENG{hide} \ENG{LMS} \ENG{UI} \ENG{controls} \\
\hline
\texttt{<\ENG{resources}>} & \ENG{Container} \ENG{for} \ENG{all} \ENG{resource} \ENG{definitions} (\ENG{SCOs}, \ENG{assets}) \\
\hline
\texttt{<\ENG{resource}>} & \ENG{Defines} \ENG{a} \ENG{SCO} \ENG{or} \ENG{asset}; \ENG{includes} \ENG{type} (\texttt{\ENG{webcontent}}), \texttt{\ENG{scormType}} (\texttt{\ENG{sco}}, \texttt{\ENG{asset}}), \ENG{and} \texttt{\ENG{href}} (\ENG{entry} \ENG{point}) \\
\hline
\texttt{<\ENG{resource} \ENG{adlcp:scormType}>} & \ENG{ADL} \ENG{extension} \ENG{indicating} \ENG{resource} \ENG{type}: \texttt{\ENG{sco}} (\ENG{communicative}) \ENG{or} \texttt{\ENG{asset}} (\ENG{non-communicative}) \\
\hline
\texttt{<\ENG{file}>} & \ENG{Declares} \ENG{files} \ENG{included} \ENG{in} \ENG{the} \ENG{content} \ENG{package} \ENG{for} \ENG{a} \ENG{resource} \\
\hline
\texttt{<\ENG{dependency}>} & \ENG{References} \ENG{shared} \ENG{resources} \ENG{used} \ENG{by} \ENG{this} \ENG{resource} \\
\hline
\end{tabular}
\end{adjustbox}
\endgroup

\end{table}

Η σημαντικότερη διαφοροποίηση σε σχέση με το \ENG{SCORM} 1.2 είναι ότι η πληροφορία \ENG{sequencing} ενσωματώνεται απευθείας στα στοιχεία \texttt{<\ENG{item}>} μέσω του \texttt{<\ENG{imsss:sequencing}>}. Έτσι, κάθε κόμβος ενός \ENG{organization} μπορεί να δηλώνει \ENG{control modes}, \ENG{limit conditions}, κανόνες \ENG{rollup}, \ENG{objectives} και άλλες συνθήκες που επηρεάζουν τη διαθεσιμότητα ή τη μετάβαση σε επόμενες δραστηριότητες. Με αυτόν τον τρόπο, σχέσεις που στο \ENG{SCORM} 1.2 αποδίδονταν αποσπασματικά μέσω απλών επεκτάσεων, όπως τα \ENG{prerequisites}, αποκτούν τώρα τυπική και αλγοριθμικά επεξεργάσιμη μορφή~\cite{IMS_SimpleSequencing, ADL_SCORM2004_SN}.

Παρά την ενίσχυση της λογικής του \ENG{manifest}, το μοντέλο πακεταρίσματος διατηρεί την αρχή του αυτοτελούς \ENG{ZIP-based} πακέτου. Το \texttt{\ENG{imsmanifest.xml}} παραμένει το κύριο έγγραφο στη ρίζα του πακέτου, ενώ το περιεχόμενο περιλαμβάνει \ENG{SCOs}, \ENG{assets}, βοηθητικά αρχεία, \ENG{schema definitions} και προαιρετικά αρχεία ελέγχου ή μεταδεδομένων. Η διαφορά είναι ότι το \ENG{SCORM} 2004 απαιτεί αυστηρότερη συμμόρφωση ως προς τις δηλώσεις \ENG{namespace}, τις αναφορές \ENG{schema} και τη σωστή χρήση των \ENG{ADL} επεκτάσεων, γεγονός που ενισχύει την τυπική συμβατότητα μεταξύ πακέτου και \ENG{LMS}~\cite{IMS_CP_Specification, ADL_SCORM2004_CAM}.

Επιπλέον, το \ENG{SCORM} 2004 αποσαφηνίζει τη σχέση ανάμεσα σε \ENG{assets}, \ENG{SCOs} και \ENG{activities}. Τα \ENG{assets} παραμένουν μη επικοινωνούντες πόροι και τα \ENG{SCOs} παραμένουν εκτελέσιμες μονάδες με πρόσβαση στο \ENG{Run-Time API}, αλλά η εισαγωγή της έννοιας του \ENG{activity} επιτρέπει στο \ENG{sequencing} υποσύστημα να χειρίζεται αφηρημένα τόσο εκτελέσιμους όσο και μη εκτελέσιμους κόμβους. Κάθε \texttt{<\ENG{item}>} του \ENG{manifest} αντιστοιχεί πλέον σε ένα \ENG{activity} του \ENG{activity tree}, η οποία μπορεί να αναφέρεται σε πόρο ή να λειτουργεί ως καθαρά δομική συσσωμάτωση. Πρόκειται για κάτι κρίσιμο, επειδή το \ENG{sequencing} δεν λειτουργεί απευθείας πάνω σε αρχεία αλλά πάνω σε δραστηριότητες και στην κατάστασή τους~\cite{ADL_SCORM2004_SN, IMS_SS_BehaviorModel}.

\subsubsection{Ενσωμάτωση \ENG{Metadata}}

Η υποστήριξη μεταδεδομένων στο \ENG{SCORM} 2004 είναι εκτενέστερη και εννοιολογικά ωριμότερη από εκείνη του \ENG{SCORM} 1.2, καθώς ενσωματώνει πλήρως την προδιαγραφή \ENG{IEEE LOM} με \ENG{SCORM}-ειδικά \ENG{application profiles}. Τα \ENG{metadata} μπορούν να εφαρμοστούν στο επίπεδο του συνολικού πακέτου, ενός συγκεκριμένου \ENG{organization}, επιμέρους \ENG{item} ή μεμονωμένου \ENG{resource}, ώστε η περιγραφή να ακολουθεί την ιεραρχία του περιεχομένου και να μην περιορίζεται μόνο σε μία συνολική καταγραφή. Το αποτέλεσμα είναι βελτιωμένη δυνατότητα καταλογογράφησης, ακριβέστερη αναζήτηση σε αποθετήρια και αποτελεσματικότερη επαναχρησιμοποίηση σε σενάρια όπου τα ίδια μαθησιακά αντικείμενα επανεντάσσονται σε διαφορετικές ακολουθίες μάθησης~\cite{IEEE_LOM_2020, ADL_SCORM2004_CAM}.

\subsection{Περιβάλλον Εκτέλεσης (\ENG{Run-Time Environment} -- \ENG{RTE})}
\label{subsec:scorm2004_rte}

Το \ENG{Run-Time Environment} του \ENG{SCORM} 2004 εξελίσσει τη λογική του \ENG{SCORM} 1.2, αλλά επιχειρεί παράλληλα να επιβάλει αυστηρότερη πειθαρχία στην ανταλλαγή δεδομένων και στον κύκλο ζωής της συνεδρίας. Η πρόθεση της προδιαγραφής είναι σαφής: το \ENG{runtime} δεν πρέπει μόνο να συλλέγει τιμές από το περιεχόμενο, αλλά να παρέχει ένα σταθερό και αυστηρά ορισμένο πλαίσιο μέσα στο οποίο η κατάσταση μάθησης μπορεί να χρησιμοποιηθεί και από το υποσύστημα \ENG{sequencing}~\cite{ADL_SCORM2004_RTE, ADL_SCORM_Users_Guide}.

\subsubsection{Προδιαγραφή \ENG{API}}

Η πιο ορατή αλλαγή στο επίπεδο διεπαφής είναι η μετονομασία του αντικειμένου \ENG{API} σε \texttt{\ENG{API}\_1484\_11}, επιλογή που αντανακλά την εναρμόνιση με το σχετικό \ENG{IEEE} πρότυπο και ταυτόχρονα διακρίνει καθαρά το \ENG{SCORM} 2004 από το \ENG{SCORM} 1.2. Παρότι ο αριθμός των βασικών μεθόδων παραμένει ουσιαστικά ο ίδιος, η ονοματολογία γίνεται απλούστερη και η σημασιολογία τους εντάσσεται σε αυστηρότερο μοντέλο κατάστασης. Ο Πίνακας~\ref{tab:scorm2004_api_methods} παρουσιάζει τις μεθόδους του \ENG{SCORM} 2004 \ENG{API}~\cite{IEEE_ECMAScript_API, ADL_SCORM2004_RTE}.

\begin{table}[htbp]
\centering
\caption{Μέθοδοι \ENG{SCORM} 2004 \ENG{API}}
\label{tab:scorm2004_api_methods}
\begingroup
\footnotesize
\setlength{\tabcolsep}{3pt}
\renewcommand{\arraystretch}{1.16}
\begin{adjustbox}{center,max width=\textwidth}
\begin{tabular}{|p{3.5cm}|p{2.5cm}|p{8cm}|}
\hline
\ENG{Method} & \ENG{Parameters} & Περιγραφή \\
\hline
\hline
\texttt{\ENG{Initialize}()} & \texttt{""} (\ENG{empty} \ENG{string}) & Ξεκινά μια \ENG{communication} \ENG{session} με το \ENG{LMS}. Πρέπει να καλείται πριν από οποιαδήποτε κλήση \texttt{\ENG{GetValue}()} ή \texttt{\ENG{SetValue}()}. Επιστρέφει \texttt{"\ENG{true}"} σε περίπτωση επιτυχίας και \texttt{"\ENG{false}"} σε περίπτωση αποτυχίας. \\
\hline
\texttt{\ENG{Terminate}()} & \texttt{""} (\ENG{empty} \ENG{string}) & Τερματίζει την \ENG{communication} \ENG{session} με το \ENG{LMS}. Πρέπει να καλείται όταν το \ENG{SCO} ολοκληρώσει την εκτέλεσή του. Επιστρέφει \texttt{"\ENG{true}"} σε περίπτωση επιτυχίας και \texttt{"\ENG{false}"} σε περίπτωση αποτυχίας. Αν έχει οριστεί το \texttt{\ENG{adl.nav.request}}, ενεργοποιεί τη διαδικασία \ENG{sequencing}. \\
\hline
\texttt{\ENG{GetValue}()} & \texttt{\ENG{element}} (\ENG{string}) & Ανακτά την τιμή ενός καθορισμένου στοιχείου του \ENG{data} \ENG{model} από το \ENG{LMS}. Επιστρέφει την τιμή ως \ENG{string}, ή κενό \ENG{string} σε περίπτωση \ENG{error}. Ελέγξτε το \texttt{\ENG{GetLastError}()} για τους αντίστοιχους \ENG{error} \ENG{codes}. \\
\hline
\texttt{\ENG{SetValue}()} & \texttt{\ENG{element}} (\ENG{string}), \texttt{\ENG{value}} (\ENG{string}) & Αποθηκεύει μια τιμή στο καθορισμένο στοιχείο του \ENG{data} \ENG{model} στο \ENG{LMS}. Επιστρέφει \texttt{"\ENG{true}"} σε περίπτωση επιτυχίας και \texttt{"\ENG{false}"} σε περίπτωση αποτυχίας. Ελέγξτε το \texttt{\ENG{GetLastError}()} για πιθανά \ENG{validation} \ENG{errors}. \\
\hline
\texttt{\ENG{Commit}()} & \texttt{""} (\ENG{empty} \ENG{string}) & Δίνει εντολή στο \ENG{LMS} να αποθηκεύσει μόνιμα όλα τα δεδομένα που ορίστηκαν μέσω της \texttt{\ENG{SetValue}()} από την τελευταία \texttt{\ENG{Initialize}()} ή \texttt{\ENG{Commit}()}. Είναι προαιρετική, αλλά συνιστάται για τη διασφάλιση της ακεραιότητας των δεδομένων. Επιστρέφει \texttt{"\ENG{true}"} σε περίπτωση επιτυχίας και \texttt{"\ENG{false}"} σε περίπτωση αποτυχίας. \\
\hline
\texttt{\ENG{GetLastError}()} & \ENG{None} & Επιστρέφει τον \ENG{error} \ENG{code} (ως \ENG{string}) που προέκυψε από την τελευταία κλήση του \ENG{API}. Ο \ENG{error} \ENG{code} \texttt{"0"} δηλώνει ότι δεν υπάρχει \ENG{error}. Ανατρέξτε στον πίνακα \ENG{SCORM} 2004 \ENG{error} \ENG{codes} για την ερμηνεία τους. \\
\hline
\texttt{\ENG{GetErrorString}()} & \texttt{\ENG{errorCode}} (\ENG{string}) & Επιστρέφει μια σύντομη κειμενική περιγραφή του καθορισμένου \ENG{error} \ENG{code}. Παρέχει τυποποιημένα \ENG{error} \ENG{messages} (π.χ. \texttt{"\ENG{Data} \ENG{Model} \ENG{Element} \ENG{Is} \ENG{Read} \ENG{Only}"}). \\
\hline
\texttt{\ENG{GetDiagnostic}()} & \texttt{\ENG{errorCode}} (\ENG{string}) & Επιστρέφει αναλυτικές, \ENG{LMS-specific} διαγνωστικές πληροφορίες για το καθορισμένο \ENG{error}. Είναι χρήσιμη για \ENG{debugging} κατά την ανάπτυξη περιεχομένου. Μπορεί να επιστρέψει κενό \ENG{string} αν δεν υπάρχουν διαθέσιμα \ENG{diagnostic} δεδομένα. \\
\hline
\end{tabular}
\end{adjustbox}
\endgroup

\end{table}

Η ροή επικοινωνίας, η οποία αποτυπώνεται στο Σχήμα~\ref{fig:scorm2004_api_flow}, αναδεικνύει τις σημαντικότερες διαφορές από το \ENG{SCORM} 1.2. Η χρήση μεθόδων όπως \texttt{\ENG{Initialize}()}, \texttt{\ENG{Terminate}()} και \texttt{\ENG{Commit}()} με συνεπή παράμετρο την κενή συμβολοσειρά, η πιο αυστηρή ακολουθία αρχικοποίησης και τερματισμού και η πληρέστερη αναφορά σφαλμάτων συγκροτούν ένα \ENG{runtime} περισσότερο τυποποιημένο και λιγότερο ανεκτικό σε αποκλίσεις υλοποίησης.

\begin{figure}[htbp]
\centering
% Figure 3.8: SCORM 2004 API Communication Flow
\begin{chapterfigurebox}
\begin{tikzpicture}[
    process/.style={rectangle, draw=blue!60, fill=blue!10, thick, minimum width=2.6cm, minimum height=0.9cm, align=center, rounded corners},
    decision/.style={diamond, draw=red!60, fill=red!10, thick, minimum width=2.1cm, minimum height=1.2cm, align=center, aspect=1.6, inner sep=1pt},
    state/.style={ellipse, draw=green!60, fill=green!10, thick, minimum width=2.4cm, minimum height=0.9cm, align=center},
    arrow/.style={->, >=stealth, thick},
    error/.style={rectangle, draw=red!80, fill=red!20, thick, minimum width=2.8cm, minimum height=0.8cm, align=center, text width=2.4cm, font=\scriptsize},
    note/.style={font=\scriptsize\itshape, text width=3cm, align=left}
]

\node[state] (launch) at (0,0) {\ENG{SCO} \ENG{Launched}};
\node[process] (discover) at (2.8,0) {\ENG{Discover}\\\ENG{API}\_1484\_11};
\node[decision] (found) at (5.7,0) {\ENG{API}\\\ENG{Found}?};
\node[process] (initialize) at (8.8,0) {\ENG{Call}\\\ENG{Initialize}("")};
\node[decision] (init_success) at (11.9,0) {\ENG{Success}?};

\node[error] (err_noapi) at (5.7,-2.3) {\ENG{Error}: \ENG{No} \ENG{API}};
\node[error] (err_init) at (11.9,-2.3) {\ENG{Error}: \ENG{Init} \ENG{Failed}\\(\ENG{Code} 102)};

\node[state, fill=green!20] (running) at (3.1,-4.7) {\ENG{Running}\\\ENG{State}};
\node[process] (exchange) at (6.3,-4.7) {\ENG{GetValue}() /\\\ENG{SetValue}() /\\\ENG{Commit}()};
\node[process] (terminate) at (9.5,-4.7) {\ENG{Call}\\\ENG{Terminate}("")};
\node[decision] (term_success) at (12.5,-4.7) {\ENG{Success}?};

\node[error] (err_term) at (12.5,-7.0) {\ENG{Error}: \ENG{Term} \ENG{Failed}\\(\ENG{Code} 111)};
\node[state, fill=gray!20] (end) at (12.5,-9.2) {\ENG{SCO} \ENG{Unloaded}};

\draw[arrow] (launch) -- (discover);
\draw[arrow] (discover) -- (found);
\draw[arrow] (found) -- node[above, font=\scriptsize] {\ENG{No}} (err_noapi);
\draw[arrow] (found) -- node[above, font=\scriptsize] {\ENG{Yes}} (initialize);
\draw[arrow] (initialize) -- (init_success);
\draw[arrow] (init_success) -- node[above, font=\scriptsize] {\ENG{No}} (err_init);
\draw[arrow] (init_success.south) |- node[pos=0.25, right, font=\scriptsize] {\ENG{Yes}} (running.north);
\draw[arrow] (running) -- (exchange);
\draw[arrow, bend left=28] (exchange.north) to node[above, font=\scriptsize] {\ENG{Loop}} (running.north);
\draw[arrow] (exchange) -- (terminate);
\draw[arrow] (terminate) -- (term_success);
\draw[arrow] (term_success) -- node[right, font=\scriptsize] {\ENG{No}} (err_term);
\draw[arrow] (term_success) -- node[right, font=\scriptsize] {\ENG{Yes}} (end);

\node[note] at (2.1,-7.2) {\ENG{Data model} \ENG{operations} \ENG{are} \ENG{valid} \ENG{only} \ENG{while} \ENG{the} \ENG{session} \ENG{is} \ENG{running}.};
\node[note] at (8.8,-7.2) {\ENG{The} \ENG{API} \ENG{enforces} \ENG{the} \ENG{state} \ENG{transition} \ENG{order}.};

\end{tikzpicture}
\end{chapterfigurebox}

\caption{Ροή Επικοινωνίας \ENG{SCORM} 2004 \ENG{API}}
\label{fig:scorm2004_api_flow}
\end{figure}

Η διαδικασία εντοπισμού του \ENG{API} παραμένει ιεραρχική, αλλά πλέον στοχεύει το αντικείμενο \texttt{\ENG{API}\_1484\_11}. Το \ENG{SCO} ελέγχει διαδοχικά το τρέχον \ENG{window}, τα γονικά παράθυρα και, εφόσον απαιτείται, την αλυσίδα \ENG{opener windows}, συνήθως με πρακτικό περιορισμό σε περιορισμένο βάθος αναζήτησης ώστε να αποφεύγονται ατέρμονοι βρόχοι. Η επιτυχής ανεύρεση του αντικειμένου δεν αρκεί από μόνη της· το περιεχόμενο πρέπει να καλέσει \texttt{\ENG{Initialize}("")} πριν από οποιαδήποτε πρόσβαση στο μοντέλο δεδομένων, διότι μόνο τότε θεωρείται ότι έχει δημιουργηθεί έγκυρη συνεδρία επικοινωνίας~\cite{ADL_SCORM2004_RTE, Rustici_SCORM_API_Guide}.

Η αυστηρότητα του \ENG{SCORM} 2004 γίνεται ιδιαίτερα εμφανής στο μοντέλο κατάστασης του \ENG{API}. Πριν από την αρχικοποίηση το \ENG{runtime} βρίσκεται στην κατάσταση \ENG{Not Initialized}, μετά από επιτυχή \texttt{\ENG{Initialize}()} στην κατάσταση \ENG{Running}, και μετά το \texttt{\ENG{Terminate}()} στην κατάσταση \ENG{Terminated}. Η προδιαγραφή ορίζει ότι κλήσεις όπως \texttt{\ENG{GetValue}()} ή \texttt{\ENG{SetValue}()} εκτός της ενεργής κατάστασης οδηγούν σε συγκεκριμένους κωδικούς σφάλματος, όπως οι 122 και 123 για ανεπίτρεπτες αναγνώσεις ή οι 132 και 133 για ανεπίτρεπτες εγγραφές. Το στοιχείο αυτό καθιστά το \ENG{SCORM} 2004 περισσότερο προβλέψιμο και επιτρέπει σαφέστερη διάγνωση προβλημάτων \ENG{runtime}~\cite{ADL_SCORM2004_RTE, SCORM_Error_Reference}.

\subsubsection{Μοντέλο δεδομένων}

Το μοντέλο δεδομένων του \ENG{SCORM} 2004 επεκτείνει ουσιαστικά το αντίστοιχο του \ENG{SCORM} 1.2, όχι μόνο προσθέτοντας νέα στοιχεία αλλά και διαχωρίζοντας έννοιες που προηγουμένως συγχέονταν. Το αποτέλεσμα είναι λεπτομερέστερη παρακολούθηση, καλύτερη υποστήριξη παιδαγωγικών στόχων και στενότερη διασύνδεση με τους μηχανισμούς \ENG{sequencing}. Η ιεραρχία του μοντέλου παρουσιάζεται στο Σχήμα~\ref{fig:data_model_hierarchy}, όπου αποτυπώνεται η εσωτερική δομή του \texttt{\ENG{cmi}.*} \ENG{namespace}~\cite{ADL_SCORM2004_RTE, IEEE_LOM_Data_Model}.

\begin{figure}[htbp]
\centering
% Figure 3.11: SCORM 2004 Data Model Hierarchy
\begin{chapterfigurebox}
\begin{tikzpicture}[
    every node/.style={font=\small},
    arrow/.style={->, >=latex, thick},
    namespace/.style={rectangle, draw=blue!80, fill=blue!20, thick, minimum width=4.6cm, minimum height=0.95cm, align=center, rounded corners},
    category/.style={rectangle, draw=green!70, fill=green!15, thick, rounded corners, text width=3.2cm, minimum height=2.55cm, align=left, font=\scriptsize},
    navbox/.style={rectangle, draw=purple!80, fill=purple!20, thick, rounded corners, text width=10.4cm, minimum height=1.45cm, align=left, font=\scriptsize}
]

\node[namespace] (cmi) at (0,0) {\texttt{\ENG{cmi}} \ENG{namespace}};

\node[category] (core_info) at (-6.0,-2.9) {\textbf{\ENG{Core Info}}\\\texttt{\ENG{learner}\_\ENG{id}}\\\texttt{\ENG{learner}\_\ENG{name}}\\\texttt{\ENG{location}}\\\texttt{\ENG{mode}}};
\node[category] (status) at (-2.0,-2.9) {\textbf{\ENG{Status}}\\\texttt{\ENG{completion}\_\ENG{status}}\\\texttt{\ENG{success}\_\ENG{status}}\\\texttt{\ENG{progress}\_\ENG{measure}}};
\node[category] (score) at (2.0,-2.9) {\textbf{\ENG{Score}}\\\texttt{\ENG{score.scaled}}\\\texttt{\ENG{score.raw}}\\\texttt{\ENG{score.min/max}}};
\node[category] (time) at (6.0,-2.9) {\textbf{\ENG{Time}}\\\texttt{\ENG{session}\_\ENG{time}}\\\texttt{\ENG{total}\_\ENG{time}}\\\texttt{\ENG{max}\_\ENG{time}\_\ENG{allowed}}};

\node[category] (objectives) at (-6.0,-6.6) {\textbf{\ENG{Objectives}}\\\texttt{\ENG{objectives.n.id}}\\\texttt{\ENG{objectives.n.score}}\\\texttt{\ENG{objectives.n.}}\\\texttt{\ENG{success}\_\ENG{status}}};
\node[category] (interactions) at (-2.0,-6.6) {\textbf{\ENG{Interactions}}\\\texttt{\ENG{interactions.n.}}\\\texttt{\ENG{id}}\\\texttt{\ENG{interactions.n.}}\\\texttt{\ENG{type}}\\\texttt{\ENG{interactions.n.}}\\\texttt{\ENG{learner}\_\ENG{response}}};
\node[category] (comments) at (2.0,-6.6) {\textbf{\ENG{Comments}}\\\texttt{\ENG{comments}\_\ENG{from}.}\\\texttt{\ENG{learner}}\\\texttt{\ENG{comments}\_\ENG{from}.}\\\texttt{\ENG{lms}}};
\node[category] (other) at (6.0,-6.6) {\textbf{\ENG{Other}}\\\texttt{\ENG{suspend}\_\ENG{data}}\\\texttt{\ENG{launch}\_\ENG{data}}\\\texttt{\ENG{learner}\_\ENG{preference}}};

\node[
    draw=gray!60,
    dashed,
    rounded corners,
    inner sep=0.60cm,
    fit=(core_info) (other),
    label={[font=\scriptsize\bfseries, anchor=north west]north west:\ENG{Representative} \ENG{data} \ENG{categories}}
] (cmi_group) {};

\draw[arrow] (cmi) -- (cmi_group.north);

\node[navbox, below=1.7cm of cmi_group] (adlnav) {\texttt{\ENG{adl.nav}} \ENG{namespace}: \texttt{\ENG{request}}, \texttt{\ENG{request}\_\ENG{valid.continue}}\\\texttt{\ENG{request}\_\ENG{valid.previous}}, \texttt{\ENG{request}\_\ENG{valid.choice}}};
\draw[arrow] (cmi_group.south) -- (adlnav.north);

\end{tikzpicture}
\end{chapterfigurebox}

\caption{Ιεραρχία Μοντέλου Δεδομένων \ENG{SCORM} 2004}
\label{fig:data_model_hierarchy}
\end{figure}

Ο Πίνακας~\ref{tab:cmi_data_model_2004} παρουσιάζει αναλυτικά τα στοιχεία του \ENG{SCORM} 2004 \ENG{CMI} μοντέλου. Η αξία του εκτεταμένου αυτού μοντέλου έγκειται στο ότι παρέχει όχι απλώς περισσότερα πεδία, αλλά μια σαφέστερη διάκριση ανάμεσα σε πρόοδο, ολοκλήρωση, επιτυχία, επίδοση, στόχους και σχόλια.

\begingroup
\small
\setlength{\tabcolsep}{3pt}
\renewcommand{\arraystretch}{1.14}
\setlength\LTleft{0pt plus 1fil}
\setlength\LTright{0pt plus 1fil}
\begin{longtable}{|>{\ttfamily\raggedright\arraybackslash}p{6.8cm}|>{\centering\arraybackslash}p{1.0cm}|>{\raggedright\arraybackslash}p{4.7cm}|}
\caption{Στοιχεία \ENG{SCORM} 2004 \ENG{CMI} Μοντέλου Δεδομένων}
\label{tab:cmi_data_model_2004}\\
\hline
\normalfont\ENG{Data} \ENG{Model} \ENG{Element} & \ENG{Access} & \ENG{Description} \\
\hline
\hline
\endfirsthead
\multicolumn{3}{c}{\tablename\ \thetable\ -- συνέχεια} \\
\hline
\normalfont\ENG{Data} \ENG{Model} \ENG{Element} & \ENG{Access} & \ENG{Description} \\
\hline
\hline
\endhead
\hline
\multicolumn{3}{r}{\footnotesize Συνεχίζεται στην επόμενη σελίδα} \\
\endfoot
\hline
\endlastfoot
\multicolumn{3}{|c|}{\ENG{Core} \ENG{Learner} \ENG{and} \ENG{Session} \ENG{Information}} \\
\hline
\texttt{\ENG{cmi}.\_\ENG{version}} & \ENG{RO} & \ENG{Data} \ENG{model} \ENG{version} (\ENG{e.g}., "1.0" \ENG{for} \ENG{SCORM} 2004 3\ENG{rd} \ENG{Ed}.) \\
\hline
\texttt{\ENG{cmi.learner}\_\ENG{id}} & \ENG{RO} & \ENG{Unique} \ENG{identifier} \ENG{for} \ENG{the} \ENG{learner} \\
\hline
\texttt{\ENG{cmi.learner}\_\ENG{name}} & \ENG{RO} & \ENG{Learner}'\ENG{s} \ENG{name} (\ENG{localized} \ENG{string}) \\
\hline
\texttt{\ENG{cmi.location}} & \ENG{RW} & \ENG{Bookmark/location} \ENG{within} \ENG{SCO} (\ENG{up} \ENG{to} 1000 \ENG{chars}) \\
\hline
\texttt{\ENG{cmi.credit}} & \ENG{RO} & \ENG{Indicates} \ENG{credit} \ENG{mode}: \texttt{\ENG{credit}}, \texttt{\ENG{no-credit}} \\
\hline
\texttt{\ENG{cmi.mode}} & \ENG{RO} & \ENG{Presentation} \ENG{mode}: \texttt{\ENG{browse}}, \texttt{\ENG{normal}}, \texttt{\ENG{review}} \\
\hline
\texttt{\ENG{cmi.entry}} & \ENG{RO} & \ENG{Entry} \ENG{point}: \texttt{\ENG{ab-initio}} (\ENG{first} \ENG{time}), \texttt{\ENG{resume}}, \texttt{""} \\
\hline
\texttt{\ENG{cmi.exit}} & \ENG{WO} & \ENG{Exit} \ENG{reason}: \texttt{\ENG{timeout}}, \texttt{\ENG{suspend}}, \texttt{\ENG{logout}}, \texttt{\ENG{normal}}, \texttt{""} \\
\hline
\hline
\multicolumn{3}{|c|}{\ENG{Completion} \ENG{and} \ENG{Success} \ENG{Status}} \\
\hline
\texttt{\ENG{cmi.completion}\_\ENG{status}} & \ENG{RW} & \ENG{Completion}: \texttt{\ENG{completed}}, \texttt{\ENG{incomplete}}, \texttt{\ENG{not} \ENG{attempted}}, \texttt{\ENG{unknown}} \\
\hline
\texttt{\ENG{cmi.completion}\_\ENG{threshold}} & \ENG{RO} & \ENG{Threshold} \ENG{for} \ENG{automatic} \ENG{completion} (\ENG{real} [0..1]) \\
\hline
\texttt{\ENG{cmi.success}\_\ENG{status}} & \ENG{RW} & \ENG{Mastery}: \texttt{\ENG{passed}}, \texttt{\ENG{failed}}, \texttt{\ENG{unknown}} \\
\hline
\texttt{\ENG{cmi.progress}\_\ENG{measure}} & \ENG{RW} & \ENG{Progress} \ENG{percentage} (\ENG{real} [0..1]) \\
\hline
\hline
\multicolumn{3}{|c|}{\ENG{Score}} \\
\hline
\texttt{\ENG{cmi.score.scaled}} & \ENG{RW} & \ENG{Normalized} \ENG{score} (\ENG{real} [-1..1]) \\
\hline
\texttt{\ENG{cmi.score.raw}} & \ENG{RW} & \ENG{Raw} \ENG{score} (\ENG{real}, \ENG{unbounded}) \\
\hline
\texttt{\ENG{cmi.score.min}} & \ENG{RW} & \ENG{Minimum} \ENG{possible} \ENG{raw} \ENG{score} \\
\hline
\texttt{\ENG{cmi.score.max}} & \ENG{RW} & \ENG{Maximum} \ENG{possible} \ENG{raw} \ENG{score} \\
\hline
\texttt{\ENG{cmi.scaled}\_\ENG{passing}\_\ENG{score}} & \ENG{RO} & \ENG{Normalized} \ENG{passing} \ENG{score} (\ENG{real} [-1..1]) \\
\hline
\hline
\multicolumn{3}{|c|}{\ENG{Time} \ENG{Tracking}} \\
\hline
\texttt{\ENG{cmi.session}\_\ENG{time}} & \ENG{WO} & \ENG{Session} \ENG{duration} (\ENG{timeinterval}, \ENG{e.g}., \texttt{\ENG{PT1H30M15S}}) \\
\hline
\texttt{\ENG{cmi.total}\_\ENG{time}} & \ENG{RO} & \ENG{Cumulative} \ENG{time} \ENG{across} \ENG{all} \ENG{attempts} (\ENG{timeinterval}) \\
\hline
\texttt{\ENG{cmi.max}\_\ENG{time}\_\ENG{allowed}} & \ENG{RO} & \ENG{Maximum} \ENG{allowed} \ENG{time} \ENG{for} \ENG{SCO} (\ENG{timeinterval}) \\
\hline
\texttt{\ENG{cmi.time}\_\ENG{limit}\_\ENG{action}} & \ENG{RO} & \ENG{Action} \ENG{on} \ENG{time} \ENG{limit}: \texttt{\ENG{exit},\ENG{message}}, \texttt{\ENG{continue},\ENG{message}}, \texttt{\ENG{exit},\ENG{no} \ENG{message}}, \texttt{\ENG{continue},\ENG{no} \ENG{message}} \\
\hline
\hline
\multicolumn{3}{|c|}{\ENG{Suspend} \ENG{Data} \ENG{and} \ENG{Launch} \ENG{Data}} \\
\hline
\texttt{\ENG{cmi.suspend}\_\ENG{data}} & \ENG{RW} & \ENG{SCO} \ENG{state} \ENG{data} (4000 \ENG{chars} \ENG{SCORM} 2004 2\ENG{nd} \ENG{Ed}.; 64000 \ENG{chars} 3\ENG{rd/4th} \ENG{Ed}.) \\
\hline
\texttt{\ENG{cmi.launch}\_\ENG{data}} & \ENG{RO} & \ENG{Data} \ENG{from} \ENG{manifest} \texttt{\ENG{dataFromLMS}} \ENG{element} (\ENG{up} \ENG{to} 4000 \ENG{chars}) \\
\hline
\hline
\multicolumn{3}{|c|}{\ENG{Comments}} \\
\hline
\texttt{\ENG{cmi.comments}\_\ENG{from}\_\ENG{learner}.\_\ENG{count}} & \ENG{RO} & \ENG{Number} \ENG{of} \ENG{learner} \ENG{comments} \\
\hline
\texttt{\ENG{cmi.comments}\_\ENG{from}\_\ENG{learner.n.comment}} & \ENG{RW} & \ENG{Comment} \ENG{text} (\ENG{localized} \ENG{string}, \ENG{up} \ENG{to} 4000 \ENG{chars}) \\
\hline
\texttt{\ENG{cmi.comments}\_\ENG{from}\_\ENG{learner.n.location}} & \ENG{RW} & \ENG{Location} \ENG{reference} \ENG{for} \ENG{comment} (\ENG{up} \ENG{to} 250 \ENG{chars}) \\
\hline
\texttt{\ENG{cmi.comments}\_\ENG{from}\_\ENG{learner.n.timestamp}} & \ENG{RW} & \ENG{Timestamp} (\ENG{ISO} 8601 \ENG{format}) \\
\hline
\texttt{\ENG{cmi.comments}\_\ENG{from}\_\ENG{lms}.\_\ENG{count}} & \ENG{RO} & \ENG{Number} \ENG{of} \ENG{LMS} \ENG{comments} \\
\hline
\texttt{\ENG{cmi.comments}\_\ENG{from}\_\ENG{lms.n.comment}} & \ENG{RO} & \ENG{LMS} \ENG{comment} \ENG{text} (\ENG{localized} \ENG{string}, \ENG{up} \ENG{to} 4000 \ENG{chars}) \\
\hline
\texttt{\ENG{cmi.comments}\_\ENG{from}\_\ENG{lms.n.location}} & \ENG{RO} & \ENG{Location} \ENG{reference} \ENG{for} \ENG{LMS} \ENG{comment} \\
\hline
\texttt{\ENG{cmi.comments}\_\ENG{from}\_\ENG{lms.n.timestamp}} & \ENG{RO} & \ENG{Timestamp} \ENG{for} \ENG{LMS} \ENG{comment} \\
\hline
\hline
\multicolumn{3}{|c|}{\ENG{Objectives}} \\
\hline
\texttt{\ENG{cmi.objectives}.\_\ENG{count}} & \ENG{RO} & \ENG{Number} \ENG{of} \ENG{objectives} \\
\hline
\texttt{\ENG{cmi.objectives.n.id}} & \ENG{RW} & \ENG{Objective} \ENG{identifier} (\ENG{up} \ENG{to} 4000 \ENG{chars}) \\
\hline
\texttt{\ENG{cmi.objectives.n.score.scaled}} & \ENG{RW} & \ENG{Objective} \ENG{score} (\ENG{real} [-1..1]) \\
\hline
\texttt{\ENG{cmi.objectives.n.score.raw}} & \ENG{RW} & \ENG{Objective} \ENG{raw} \ENG{score} \\
\hline
\texttt{\ENG{cmi.objectives.n.score.min}} & \ENG{RW} & \ENG{Objective} \ENG{min} \ENG{score} \\
\hline
\texttt{\ENG{cmi.objectives.n.score.max}} & \ENG{RW} & \ENG{Objective} \ENG{max} \ENG{score} \\
\hline
\texttt{\ENG{cmi.objectives.n.success}\_\ENG{status}} & \ENG{RW} & \ENG{Objective} \ENG{mastery}: \texttt{\ENG{passed}}, \texttt{\ENG{failed}}, \texttt{\ENG{unknown}} \\
\hline
\texttt{\ENG{cmi.objectives.n.completion}\_\ENG{status}} & \ENG{RW} & \ENG{Objective} \ENG{completion}: \texttt{\ENG{completed}}, \texttt{\ENG{incomplete}}, \texttt{\ENG{not} \ENG{attempted}}, \texttt{\ENG{unknown}} \\
\hline
\texttt{\ENG{cmi.objectives.n.progress}\_\ENG{measure}} & \ENG{RW} & \ENG{Objective} \ENG{progress} (\ENG{real} [0..1]) \\
\hline
\texttt{\ENG{cmi.objectives.n.description}} & \ENG{RW} & \ENG{Objective} \ENG{description} (\ENG{localized} \ENG{string}, \ENG{up} \ENG{to} 250 \ENG{chars}) \\
\hline
\hline
\multicolumn{3}{|c|}{\ENG{Interactions}} \\
\hline
\texttt{\ENG{cmi.interactions}.\_\ENG{count}} & \ENG{RO} & \ENG{Number} \ENG{of} \ENG{interactions} \\
\hline
\texttt{\ENG{cmi.interactions.n.id}} & \ENG{RW} & \ENG{Interaction} \ENG{identifier} (\ENG{up} \ENG{to} 4000 \ENG{chars}) \\
\hline
\texttt{\ENG{cmi.interactions.n.type}} & \ENG{RW} & \ENG{Interaction} \ENG{type}: \texttt{\ENG{true-false}}, \texttt{\ENG{choice}}, \texttt{\ENG{fill-in}}, \texttt{\ENG{long-fill-in}}, \texttt{\ENG{matching}}, \texttt{\ENG{performance}}, \texttt{\ENG{sequencing}}, \texttt{\ENG{likert}}, \texttt{\ENG{numeric}}, \texttt{\ENG{other}} \\
\hline
\texttt{\ENG{cmi.interactions.n.timestamp}} & \ENG{RW} & \ENG{Interaction} \ENG{timestamp} (\ENG{ISO} 8601) \\
\hline
\texttt{\ENG{cmi.interactions.n.weighting}} & \ENG{RW} & \ENG{Interaction} \ENG{weight} (\ENG{real}, \ENG{unbounded}) \\
\hline
\texttt{\ENG{cmi.interactions.n.learner}\_\ENG{response}} & \ENG{RW} & \ENG{Learner} \ENG{response} (\ENG{format} \ENG{depends} \ENG{on} \ENG{type}) \\
\hline
\texttt{\ENG{cmi.interactions.n.result}} & \ENG{RW} & \ENG{Interaction} \ENG{result}: \texttt{\ENG{correct}}, \texttt{\ENG{incorrect}}, \texttt{\ENG{unanticipated}}, \texttt{\ENG{neutral}}, \ENG{or} \ENG{real} \ENG{number} \\
\hline
\texttt{\ENG{cmi.interactions.n.latency}} & \ENG{RW} & \ENG{Response} \ENG{latency} (\ENG{timeinterval}) \\
\hline
\texttt{\ENG{cmi.interactions.n.description}} & \ENG{RW} & \ENG{Interaction} \ENG{description} (\ENG{localized} \ENG{string}, \ENG{up} \ENG{to} 250 \ENG{chars}) \\
\hline
\texttt{\ENG{cmi.interactions.n.objectives}.\_\ENG{count}} & \ENG{RO} & \ENG{Number} \ENG{of} \ENG{objectives} \ENG{for} \ENG{interaction} \\
\hline
\texttt{\ENG{cmi.interactions.n.objectives.n.id}} & \ENG{RW} & \ENG{Objective} \ENG{ID} \ENG{associated} \ENG{with} \ENG{interaction} \\
\hline
\texttt{\ENG{cmi.interactions.n.correct}\_\ENG{responses}.\_\ENG{count}} & \ENG{RO} & \ENG{Number} \ENG{of} \ENG{correct} \ENG{response} \ENG{patterns} \\
\hline
\texttt{\ENG{cmi.interactions.n.correct}\_\ENG{responses.n.pattern}} & \ENG{RW} & \ENG{Correct} \ENG{response} \ENG{pattern} (\ENG{format} \ENG{depends} \ENG{on} \ENG{type}) \\
\hline
\hline
\multicolumn{3}{|c|}{\ENG{Learner} \ENG{Preferences}} \\
\hline
\texttt{\ENG{cmi.learner}\_\ENG{preference.audio}\_\ENG{level}} & \ENG{RW} & \ENG{Audio} \ENG{level} \ENG{preference} (\ENG{real} [0..*]) \\
\hline
\texttt{\ENG{cmi.learner}\_\ENG{preference.language}} & \ENG{RW} & \ENG{Preferred} \ENG{language} (\ENG{language} \ENG{code}, \ENG{up} \ENG{to} 250 \ENG{chars}) \\
\hline
\texttt{\ENG{cmi.learner}\_\ENG{preference.delivery}\_\ENG{speed}} & \ENG{RW} & \ENG{Delivery} \ENG{speed} \ENG{preference} (\ENG{real} [0..*]) \\
\hline
\texttt{\ENG{cmi.learner}\_\ENG{preference.audio}\_\ENG{captioning}} & \ENG{RW} & \ENG{Captioning} \ENG{preference}: \texttt{-1} (\ENG{off}), \texttt{0} (\ENG{no} \ENG{change}), \texttt{1} (\ENG{on}) \\
\hline
\hline
\multicolumn{3}{|c|}{\ENG{ADL} \ENG{Navigation} (\ENG{adl.nav} \ENG{namespace})} \\
\hline
\texttt{\ENG{adl.nav.request}} & \ENG{RW} & \ENG{Navigation} \ENG{request}: \texttt{\ENG{continue}}, \texttt{\ENG{previous}}, \texttt{\ENG{choice}}, \texttt{\ENG{exit}}, \texttt{\ENG{exitAll}}, \texttt{\ENG{abandon}}, \texttt{\ENG{abandonAll}}, \texttt{\ENG{suspendAll}}, \texttt{\_\ENG{none}\_} \\
\hline
\texttt{\ENG{adl.nav.request}\_\ENG{valid.continue}} & \ENG{RO} & \ENG{Can} \ENG{continue}? \texttt{\ENG{true}}, \texttt{\ENG{false}}, \texttt{\ENG{unknown}} \\
\hline
\texttt{\ENG{adl.nav.request}\_\ENG{valid.previous}} & \ENG{RO} & \ENG{Can} \ENG{go} \ENG{previous}? \texttt{\ENG{true}}, \texttt{\ENG{false}}, \texttt{\ENG{unknown}} \\
\hline
\texttt{\ENG{adl.nav.request}\_\ENG{valid.choice}.\{\ENG{target}=\}} & \ENG{RO} & \ENG{Can} \ENG{choose} \ENG{target}? \texttt{\ENG{true}}, \texttt{\ENG{false}}, \texttt{\ENG{unknown}} (\ENG{target} \ENG{specified} \ENG{in} \ENG{curly} \ENG{braces}) \\
\hline
\end{longtable}
\endgroup


Η σημαντικότερη ίσως βελτίωση είναι ο ρητός διαχωρισμός ανάμεσα στο \texttt{\ENG{cmi.completion}\_\ENG{status}} και στο \texttt{\ENG{cmi.success}\_\ENG{status}}. Στο \ENG{SCORM} 1.2 η κατάσταση του μαθήματος όφειλε να χωρέσει ταυτόχρονα έννοιες ολοκλήρωσης και επίτευξης, ενώ στο \ENG{SCORM} 2004 η ολοκλήρωση μιας δραστηριότητας και η επιτυχία σε αυτήν αποτελούν δύο ανεξάρτητες διαστάσεις. Η επιλογή αυτή επιτρέπει πιο ακριβή παιδαγωγικά σενάρια, όπου ο εκπαιδευόμενος μπορεί να έχει ολοκληρώσει μια δραστηριότητα χωρίς να την έχει περάσει ή, αντιστρόφως, να έχει αποδείξει επάρκεια χωρίς να έχει καταναλώσει όλο το προβλεπόμενο περιεχόμενο.

Η ίδια λογική ακρίβειας επεκτείνεται στην παρακολούθηση προόδου και βαθμολογίας. Τα στοιχεία \texttt{\ENG{cmi.progress}\_\ENG{measure}} και \texttt{\ENG{cmi.completion}\_\ENG{threshold}} επιτρέπουν την έκφραση προόδου ως πραγματικού αριθμού στο διάστημα [0..1], υπερβαίνοντας τα δυαδικά σχήματα \ENG{completed/incomplete}. Αντίστοιχα, η εισαγωγή του \texttt{\ENG{cmi.score.scaled}} και του \texttt{\ENG{cmi.scaled}\_\ENG{passing}\_\ENG{score}} επιτρέπει κανονικοποιημένη αξιολόγηση και καθιστά εφικτή τη σύγκριση αποτελεσμάτων μεταξύ ετερογενών δραστηριοτήτων. Οι τιμές \ENG{raw}, \ENG{min} και \ENG{max} εξακολουθούν να υπάρχουν, αλλά πλέον εντάσσονται σε πιο συνεκτικό μοντέλο μέτρησης.

Το υποσύστημα \ENG{interactions} ενισχύεται επίσης σημαντικά, υποστηρίζοντας επιπλέον τύπους αλληλεπίδρασης όπως \ENG{long-fill-in} και \ENG{other}, πλουσιότερες χρονικές σημάνσεις, αναλυτικότερη καταγραφή \ENG{learner response} και περισσότερα μοτίβα ορθών απαντήσεων. Παράλληλα, το όριο του \ENG{suspend data} αυξάνεται δραστικά στις νεότερες εκδόσεις του προτύπου, φθάνοντας τους 64000 χαρακτήρες στις \ENG{3rd} και \ENG{4th Edition}, γεγονός που καθιστά ρεαλιστική την αποθήκευση σύνθετης κατάστασης για μεγάλα ή προσαρμοστικά μαθήματα. Στο ίδιο πνεύμα, το μοντέλο \ENG{objectives} αποκτά χωριστή παρακολούθηση \ENG{success status}, \ENG{completion status}, \ENG{progress measure} και περιγραφικού κειμένου για κάθε στόχο, ενώ το μοντέλο σχολίων διασπάται στα \texttt{\ENG{cmi.comments}\_\ENG{from}\_\ENG{learner}} και \texttt{\ENG{cmi.comments}\_\ENG{from}\_\ENG{lms}}, επιτρέποντας δομημένα και χρονικά χαρακτηρισμένα σχόλια και από τις δύο πλευρές.

Μια επιπλέον διαφοροποίηση με άμεση επίδραση στη ροή του μαθήματος είναι το \ENG{adl.nav} \ENG{namespace}, το οποίο συνδέει το \ENG{RTE} με το υποσύστημα \ENG{sequencing}. Το στοιχείο \texttt{\ENG{adl.nav.request}} επιτρέπει στο \ENG{SCO} να δηλώσει, πριν από το \texttt{\ENG{Terminate}()}, αίτημα όπως \ENG{continue}, \ENG{previous}, \ENG{choice}, \ENG{exit}, \ENG{exitAll}, \ENG{abandon}, \ENG{abandonAll} ή \ENG{suspendAll}. Παράλληλα, τα στοιχεία \texttt{\ENG{adl.nav.request}\_\ENG{valid}.*} επιτρέπουν στον κώδικα του \ENG{SCO} να ελέγξει εκ των προτέρων αν ένα τέτοιο αίτημα είναι επί του παρόντος επιτρεπτό. Έτσι, το περιεχόμενο παύει να είναι απλός αποστολέας στοιχείων παρακολούθησης και αποκτά περιορισμένη αλλά ουσιαστική συμμετοχή στη λογική πλοήγησης~\cite{ADL_SCORM2004_SN, IMS_SS_BehaviorModel}.

\subsubsection{Διαχείριση σφαλμάτων}

Το \ENG{SCORM} 2004 υιοθετεί εκτεταμένο και περισσότερο συστηματικό σχήμα κωδικών σφαλμάτων, το οποίο διακρίνει με σαφήνεια τις γενικές αποτυχίες συνεδρίας, τα σφάλματα ορισμάτων, τις αποτυχίες ανάγνωσης και εγγραφής, καθώς και προβλήματα που σχετίζονται με τον ορισμό, την υποστήριξη ή τους περιορισμούς των στοιχείων του μοντέλου δεδομένων. Έτσι, οι κατηγορίες 1\ENG{xx} αντιστοιχούν σε γενικές εξαιρέσεις και αποτυχίες αρχικοποίησης ή τερματισμού, οι 2\ENG{xx} αφορούν ορίσματα, οι 3\ENG{xx}, 35\ENG{x} και 39\ENG{x} συνδέονται με \ENG{Get}, \ENG{Set} και \ENG{Commit}, ενώ οι 4\ENG{xx} περιγράφουν παραβιάσεις συμβολαίου στο επίπεδο του μοντέλου δεδομένων, όπως \ENG{read-only} εγγραφές, μη υλοποιημένα στοιχεία ή σφάλματα εύρους τιμών.

Η αναφορά σφαλμάτων παραμένει τριεπίπεδη μέσω των \texttt{\ENG{GetLastError}()}, \texttt{\ENG{GetErrorString}()} και \texttt{\ENG{GetDiagnostic}()}. Ο πρώτος μηχανισμός αποδίδει τον αριθμητικό κωδικό, ο δεύτερος παρέχει συνοπτική τυποποιημένη περιγραφή και ο τρίτος επιστρέφει λεπτομερέστερα διαγνωστικά στοιχεία, τα οποία είναι συχνά εξαρτημένα από τη συγκεκριμένη υλοποίηση του \ENG{LMS}. Η κλιμάκωση αυτή είναι κρίσιμη για την ανάπτυξη και τον έλεγχο \ENG{SCOs}, επειδή επιτρέπει γρήγορη αναγνώριση του γενικού τύπου σφάλματος αλλά και βαθύτερη διερεύνηση όταν απαιτείται~\cite{ADL_SCORM2004_RTE, SCORM_Error_Reference}.

\subsection{\ENG{Sequencing} και \ENG{Navigation} (\ENG{SN})}
\label{subsec:scorm2004_sn}

Η προδιαγραφή \ENG{Sequencing and Navigation} αποτελεί την ουσιαστικότερη καινοτομία του \ENG{SCORM} 2004, καθώς μεταφέρει την εκπαιδευτική ροή από ένα άτυπο επίπεδο διεπαφής του \ENG{LMS} σε ένα ρητά ορισμένο και τυποποιημένο υποσύστημα. Με αφετηρία το \ENG{IMS Simple Sequencing}, το \ENG{SCORM} 2004 επιτρέπει προαπαιτούμενα, εναλλακτικές διαδρομές, \ENG{remediation loops}, \ENG{randomization}, \ENG{test-out} σενάρια και γενικότερα προσαρμοστική πλοήγηση μέσα στο ίδιο πακέτο περιεχομένου~\cite{IMS_SimpleSequencing, ADL_SCORM2004_SN}.

\subsubsection{Βασικές αρχές \ENG{Sequencing}}

Το \ENG{sequencing} λειτουργεί πάνω σε ένα \ENG{activity tree} που προκύπτει από το \ENG{organization} του \ENG{manifest}. Κάθε \texttt{<\ENG{item}>} αντιστοιχεί σε μία \ENG{activity}, η οποία μπορεί να είναι \ENG{leaf activity} όταν αναφέρεται σε εκτελέσιμο πόρο ή \ENG{cluster activity} όταν λειτουργεί ως γονικός κόμβος που περιέχει άλλες δραστηριότητες. Η επιλογή της επόμενης \ENG{activity} δεν καθορίζεται πλέον μόνο από το αν ο χρήστης πάτησε ``επόμενο'' ή επέλεξε στοιχείο από έναν πίνακα περιεχομένων, αλλά από τον συνδυασμό αιτημάτων \ENG{navigation}, κανόνων \ENG{sequencing}, κατάστασης παρακολούθησης και περιορισμών ελέγχου ροής.

Το Σχήμα~\ref{fig:sequencing_model} παρουσιάζει τη γενική αρχιτεκτονική του \ENG{sequencing model} και δείχνει ότι η απόφαση για την επόμενη δραστηριότητα είναι προϊόν αλγοριθμικής επεξεργασίας και όχι απλής ευθύγραμμης πλοήγησης~\cite{ADL_SCORM2004_SN, IMS_SS_BehaviorModel}.

\begin{figure}[htbp]
\centering
% Figure 3.9: SCORM 2004 Sequencing Model
\begin{chapterfigurebox}
\begin{tikzpicture}[
    node distance=1.2cm and 2cm,
    component/.style={rectangle, draw=blue!60, fill=blue!10, thick, minimum width=3.2cm, minimum height=1cm, align=center, rounded corners},
    process/.style={rectangle, draw=green!60, fill=green!10, thick, minimum width=2.8cm, minimum height=0.9cm, align=center},
    data/.style={cylinder, draw=orange!60, fill=orange!10, thick, minimum width=2.5cm, minimum height=1cm, align=center, shape border rotate=90},
    decision/.style={diamond, draw=red!60, fill=red!10, thick, minimum width=2.2cm, minimum height=1.2cm, align=center, aspect=2},
    arrow/.style={->, >=stealth, thick},
    bidirectional/.style={<->, >=stealth, thick, dashed}
]

% Activity Tree
\node[data] (activity_tree) {\ENG{Activity}\\\ENG{Tree}};

% Navigation Request
\node[component, below=of activity_tree] (nav_request) {\ENG{Navigation}\\\ENG{Request}};
\draw[arrow] (nav_request) -- node[right, font=\scriptsize] {\ENG{Parse}} (activity_tree);

% Sequencing Process
\node[process, below=of nav_request] (end_current) {\ENG{End} \ENG{Current}\\\ENG{Activity}};
\draw[arrow] (nav_request) -- (end_current);

\node[process, below=of end_current] (rollup) {\ENG{Rollup}\\\ENG{Process}};
\draw[arrow] (end_current) -- node[right, font=\scriptsize] {\ENG{Transfer} \ENG{data}} (rollup);

\node[data, left=of rollup] (tracking_data) {\ENG{Activity}\\\ENG{Tracking}\\\ENG{Data}};
\draw[bidirectional] (rollup) -- (tracking_data);
\draw[bidirectional] (activity_tree.west) -| (tracking_data.north);

% Activity Selection
\node[process, below=of rollup] (select_activity) {\ENG{Select} \ENG{Next}\\\ENG{Activity}};
\draw[arrow] (rollup) -- (select_activity);
\draw[bidirectional] (select_activity.west) -- ++(-1,0) |- (activity_tree.west);

% Check Sequencing Rules
\node[decision, below=of select_activity] (check_rules) {\ENG{Rules}\\\ENG{Permit}?};
\draw[arrow] (select_activity) -- (check_rules);

\node[component, right=2.4cm of check_rules] (exit_course) {\ENG{Exit} \ENG{Course}};
\draw[arrow] (check_rules) -- node[above, font=\scriptsize] {\ENG{No}} (exit_course);

% Deliver Activity
\node[component, below=of check_rules] (deliver) {\ENG{Deliver}\\\ENG{Activity}};
\draw[arrow] (check_rules) -- node[left, font=\scriptsize] {\ENG{Yes}} (deliver);

% SCO Execution
\node[component, below=of deliver, fill=yellow!10, draw=yellow!60] (sco_exec) {\ENG{SCO} \ENG{Execution}\\(\ENG{RTE})};
\draw[arrow] (deliver) -- (sco_exec);

% Loop back
\draw[arrow, blue!50, bend right=50] (sco_exec.east) to node[right, font=\scriptsize, align=left] {\ENG{Terminate} +\\\ENG{Nav} \ENG{Request}} (nav_request.east);

% Sequencing Rules Box
\node[component, right=1.5cm of select_activity, fill=purple!10, draw=purple!60, text width=2.8cm] (seq_rules) {\ENG{Sequencing} \ENG{Rules}:\\- \ENG{Preconditions}\\- \ENG{Post-conditions}\\- \ENG{Rollup}\\- \ENG{Limits}};
\draw[bidirectional] (select_activity) -- (seq_rules);
\draw[bidirectional] (rollup) -| (seq_rules);

\end{tikzpicture}
\end{chapterfigurebox}

\caption{Μοντέλο \ENG{Sequencing} του \ENG{SCORM} 2004}
\label{fig:sequencing_model}
\end{figure}

Για να καταστεί αυτή η επεξεργασία δυνατή, κάθε \ENG{activity} διατηρεί δικό της σύνολο δεδομένων \ENG{tracking}. Στις \ENG{leaf activities} τα δεδομένα αυτά προκύπτουν κυρίως από τα στοιχεία του \ENG{RTE}, όπως \ENG{completion status}, \ENG{success status}, \ENG{progress measure} και \ENG{score}, ενώ στις \ENG{cluster activities} η κατάσταση προκύπτει από διαδικασίες \ENG{rollup} που συγκεντρώνουν πληροφορία από τα παιδιά. Η έννοια του \ENG{activity tracking} είναι επομένως ο συνδετικός κρίκος ανάμεσα στο \ENG{runtime} και στο \ENG{sequencing}: χωρίς τυποποιημένη κατάσταση δραστηριότητας δεν είναι δυνατό να εφαρμοστούν αξιόπιστα οι κανόνες ροής.

Στην ίδια λογική, ο ορισμός \ENG{sequencing} που δηλώνεται σε κάθε \texttt{<\ENG{item}>} μέσω του \texttt{<\ENG{imsss:sequencing}>} δεν πρέπει να νοηθεί ως απλή συλλογή σημαίων ρύθμισης. Αντίθετα, αποτελεί δηλωτικό μοντέλο που καθορίζει τους επιτρεπτούς τύπους πλοήγησης, τα όρια προσπαθειών, τις συνθήκες διαθεσιμότητας, τις στρατηγικές \ENG{rollup}, τον τρόπο επιλογής παιδιών και τη χαρτογράφηση στόχων. Ο Πίνακας~\ref{tab:sequencing_rules} συνοψίζει αυτές τις βασικές κατηγορίες ελέγχων~\cite{IMS_SimpleSequencing, ADL_SCORM2004_SN}.

\begin{table}[htbp]
\centering
\caption{Κανόνες και Έλεγχοι \ENG{SCORM} 2004 \ENG{Sequencing}}
\label{tab:sequencing_rules}
\begingroup
\footnotesize
\setlength{\tabcolsep}{3pt}
\renewcommand{\arraystretch}{1.16}
\begin{adjustbox}{center,max width=\textwidth,max totalheight=0.86\textheight}
\begin{tabular}{|p{4.5cm}|p{9.5cm}|}
\hline
\ENG{Sequencing} \ENG{Category} & \ENG{Description} \ENG{and} \ENG{Key} \ENG{Elements} \\
\hline
\hline
\ENG{Sequencing} \ENG{Control} \ENG{Modes} & \ENG{Determine} \ENG{navigation} \ENG{permissions} \ENG{for} \ENG{cluster} \ENG{children}. \texttt{\ENG{choice}} (\ENG{default}: \ENG{true}) \ENG{enables} \ENG{TOC} \ENG{navigation}; \texttt{\ENG{flow}} (\ENG{default}: \ENG{false}) \ENG{enables} \ENG{previous/next}; \texttt{\ENG{forwardOnly}} (\ENG{default}: \ENG{false}) \ENG{restricts} \ENG{to} \ENG{forward-only}; \texttt{\ENG{choiceExit}} (\ENG{default}: \ENG{true}) \ENG{allows} \ENG{exiting} \ENG{via} \ENG{TOC}. \\
\hline
\ENG{Constrain} \ENG{Choice} \ENG{Controls} & \ENG{Restrict} \ENG{TOC-based} \ENG{navigation}. \texttt{\ENG{constrainChoice}} \ENG{limits} \ENG{choices} \ENG{to} \ENG{logically} "\ENG{next}" \ENG{activity}; \texttt{\ENG{preventActivation}} \ENG{prevents} \ENG{deep-linking} \ENG{into} \ENG{children}. \\
\hline
\ENG{Sequencing} \ENG{Rules} & \ENG{If-then} \ENG{conditions} \ENG{evaluated} \ENG{against} \ENG{tracking} \ENG{data}. \ENG{Precondition} \ENG{rules} (\ENG{skip}, \ENG{disabled}, \ENG{hiddenFromChoice}, \ENG{stopForwardTraversal}) \ENG{evaluated} \ENG{before} \ENG{delivery}. \ENG{Post-condition} \ENG{rules} (\ENG{exitParent}, \ENG{exitAll}, \ENG{retry}, \ENG{retryAll}, \ENG{continue}, \ENG{previous}) \ENG{evaluated} \ENG{after} \ENG{termination}. \ENG{Exit} \ENG{rules} (\ENG{exit}) \ENG{evaluated} \ENG{when} \ENG{child} \ENG{causes} \ENG{parent} \ENG{exit}. \\
\hline
\ENG{Limit} \ENG{Conditions} & \ENG{Define} \ENG{attempt} \ENG{and} \ENG{duration} \ENG{limits}. \texttt{\ENG{attemptLimit}} \ENG{specifies} \ENG{maximum} \ENG{attempts}; \texttt{\ENG{attemptAbsoluteDurationLimit}} \ENG{specifies} \ENG{maximum} \ENG{time} (\ENG{rarely} \ENG{supported} \ENG{by} \ENG{LMSs}). \\
\hline
\ENG{Rollup} \ENG{Rules} & \ENG{Define} \ENG{how} \ENG{child} \ENG{status} \ENG{aggregates} \ENG{to} \ENG{parent}. \ENG{Rules} \ENG{specify} \texttt{\ENG{childActivitySet}} (\ENG{all}, \ENG{any}, \ENG{none}, \ENG{atLeastCount}, \ENG{atLeastPercent}), \texttt{\ENG{conditions}} (\ENG{satisfied}, \ENG{completed}, \ENG{etc}.), \ENG{and} \texttt{\ENG{actions}} (\ENG{satisfied}, \ENG{notSatisfied}, \ENG{completed}, \ENG{incomplete}). \\
\hline
\ENG{Rollup} \ENG{Controls} & \ENG{Broad} \ENG{control} \ENG{over} \ENG{rollup} \ENG{participation}. \texttt{\ENG{rollupObjectiveSatisfied}} \ENG{and} \texttt{\ENG{rollupProgressCompletion}} \ENG{determine} \ENG{whether} \ENG{activity} \ENG{contributes} \ENG{to} \ENG{parent} \ENG{rollup}. \texttt{\ENG{rollupObjectiveMeasureWeight}} \ENG{weights} \ENG{activity}'\ENG{s} \ENG{measure} \ENG{in} \ENG{parent} \ENG{measure} \ENG{calculation}. \\
\hline
\ENG{Objectives} & \ENG{Local} \ENG{and} \ENG{global} \ENG{objectives} \ENG{for} \ENG{tracking} \ENG{learning} \ENG{goals}. \ENG{Each} \ENG{objective} \ENG{has} \texttt{\ENG{id}}, \texttt{\ENG{satisfiedByMeasure}} (\ENG{use} \ENG{scaled} \ENG{score} \ENG{to} \ENG{determine} \ENG{satisfaction}), \ENG{and} \texttt{\ENG{minNormalizedMeasure}} (\ENG{passing} \ENG{threshold}). \ENG{Objective} \ENG{maps} \ENG{link} \ENG{local} \ENG{objectives} \ENG{to} \ENG{global} \ENG{shared} \ENG{objectives}. \\
\hline
\ENG{Selection} \ENG{Controls} & \ENG{Random} \ENG{selection} \ENG{of} \ENG{child} \ENG{subset}. \texttt{\ENG{selectionTiming}} (\ENG{never}, \ENG{once}, \ENG{onEachNewAttempt}) \ENG{and} \texttt{\ENG{selectionCount}} \ENG{specify} \ENG{when} \ENG{and} \ENG{how} \ENG{many} \ENG{children} \ENG{to} \ENG{randomly} \ENG{select}. \\
\hline
\ENG{Randomization} \ENG{Controls} & \ENG{Shuffle} \ENG{child} \ENG{activity} \ENG{order}. \texttt{\ENG{randomizationTiming}} (\ENG{never}, \ENG{once}, \ENG{onEachNewAttempt}) \ENG{and} \texttt{\ENG{randomizeChildren}} (\ENG{boolean}) \ENG{control} \ENG{randomization}. \\
\hline
\ENG{Delivery} \ENG{Controls} & \ENG{Control} \ENG{tracking} \ENG{and} \ENG{automatic} \ENG{status} \ENG{setting}. \texttt{\ENG{tracked}} (\ENG{default}: \ENG{true}) \ENG{enables/disables} \ENG{tracking}; \texttt{\ENG{completionSetByContent}} \ENG{and} \texttt{\ENG{objectiveSetByContent}} \ENG{determine} \ENG{whether} \ENG{SCO} \ENG{or} \ENG{sequencer} \ENG{sets} \ENG{status}. \\
\hline
\ENG{Completion} \ENG{Threshold} & (4\ENG{th} \ENG{Ed}.) \ENG{Use} \ENG{progress} \ENG{measure} \ENG{to} \ENG{determine} \ENG{completion}. \texttt{\ENG{completedByMeasure}} \ENG{enables} \ENG{threshold-based} \ENG{completion}; \texttt{\ENG{minProgressMeasure}} \ENG{defines} \ENG{threshold}; \texttt{\ENG{progressWeight}} \ENG{weights} \ENG{activity}'\ENG{s} \ENG{progress} \ENG{in} \ENG{parent} \ENG{progress} \ENG{rollup}. \\
\hline
\ENG{Navigation} \ENG{Controls} & \ENG{Hide} \ENG{LMS} \ENG{UI} \ENG{elements}. \texttt{\ENG{hideLMSUI}} \ENG{specifies} \ENG{which} \ENG{controls} \ENG{to} \ENG{hide} (\ENG{previous}, \ENG{continue}, \ENG{exit}, \ENG{exitAll}, \ENG{abandon}, \ENG{abandonAll}, \ENG{suspendAll}). \\
\hline
\end{tabular}
\end{adjustbox}
\endgroup

\end{table}

\subsubsection{Κανόνες \ENG{Sequencing}}

Οι κανόνες \ENG{sequencing} είναι λογικές συνθήκες τύπου \ENG{if-then} που επενεργούν πάνω στην κατάσταση των δραστηριοτήτων. Κάθε κανόνας συνδυάζει συνθήκες, λογικούς τελεστές και ενέργειες, ώστε να αποφαίνεται είτε για τη διαθεσιμότητα μιας \ENG{activity} είτε για τις ενέργειες που θα ακολουθήσουν μετά την ολοκλήρωσή της. Η προδιαγραφή διακρίνει τρεις βασικές στιγμές αξιολόγησης. Οι \ENG{precondition rules} εφαρμόζονται πριν από την παράδοση και χρησιμοποιούνται για να παρακάμπτουν, να απενεργοποιούν ή να αποκρύπτουν δραστηριότητες, υλοποιώντας στην πράξη προαπαιτούμενα και περιορισμούς προώθησης. Οι \ENG{post-condition rules} αξιολογούνται μετά τον τερματισμό μιας \ENG{activity} και μπορούν να οδηγήσουν σε ενέργειες όπως \ENG{retry}, \ENG{retryAll}, \ENG{continue}, \ENG{previous}, \ENG{exitParent} ή \ENG{exitAll}, άρα αποτελούν το βασικό εργαλείο για \ENG{remediation loops} και εξαναγκασμένες ακολουθίες. Τέλος, οι \ENG{exit rules} ενεργοποιούνται όταν μία παιδική δραστηριότητα προκαλεί αίτημα εξόδου προς το γονικό \ENG{cluster}, επιτρέποντας ελεγχόμενη διάδοση αυτών των αιτημάτων προς τα ανώτερα επίπεδα του \ENG{activity tree}.

Η πρακτική μορφή ενός τέτοιου κανόνα φαίνεται στο ακόλουθο παράδειγμα, όπου το \ENG{sequencing} χρησιμοποιείται για την υλοποίηση σχέσης \ENG{prerequisite}.

{\selectlanguage{english}\fontencoding{T1}\selectfont

\begin{verbatim}
<imsss:sequencingRule>
  <imsss:ruleConditions conditionCombination="all">
    <imsss:ruleCondition 
      referencedObjective="prereq_objective" 
      condition="satisfied"/>
  </imsss:ruleConditions>
  <imsss:ruleAction action="disabled"/>
</imsss:sequencingRule>
\end{verbatim}
}

Ο κανόνας αυτός καθιστά διαθέσιμη τη δραστηριότητα μόνο όταν το \ENG{objective} με αναγνωριστικό ``\ENG{prereq}\_\ENG{objective}'' έχει χαρακτηριστεί ως \ENG{satisfied}, αποτυπώνοντας με τυπικό τρόπο μια λογική εξάρτησης που στο \ENG{SCORM} 1.2 μπορούσε να αποδοθεί μόνο πολύ πιο αδρά~\cite{ADL_SCORM2004_SN, IMS_SS_InfoModel}.

\subsubsection{Κανόνες και διαδικασίες \ENG{Rollup}}

Το \ENG{rollup} είναι η διαδικασία με την οποία η κατάσταση των παιδικών \ENG{activities} συγκεντρώνεται ώστε να παραχθεί η κατάσταση μιας γονικής \ENG{cluster activity}. Ο μηχανισμός αυτός είναι καίριος για τη λειτουργία σύνθετων δομών μαθήματος, επειδή επιτρέπει στο σύστημα να συμπεραίνει αν μια ενότητα έχει ολοκληρωθεί, έχει ικανοποιηθεί ή έχει αποτύχει με βάση τα δεδομένα που προκύπτουν από τις επιμέρους δραστηριότητές της~\cite{ADL_SCORM2004_SN, IMS_SS_RollupModel}.

Όταν ο δημιουργός περιεχομένου δεν ορίζει ρητούς κανόνες, το \ENG{SCORM} 2004 εφαρμόζει προεπιλεγμένη συμπεριφορά \ENG{rollup}. Ένα \ENG{cluster} θεωρείται ολοκληρωμένο όταν όλα τα \ENG{tracked children} του είναι ολοκληρωμένα, θεωρείται ικανοποιημένο όταν όλα τα σχετικά παιδιά είναι \ENG{satisfied}, ενώ η αριθμητική του μέτρηση παράγεται ως σταθμισμένος μέσος όρος των αντίστοιχων μετρήσεων των παιδιών. Οι προεπιλεγμένοι αυτοί κανόνες επιτρέπουν λειτουργία χωρίς πρόσθετη ρύθμιση, αλλά συχνά δεν επαρκούν για πιο σύνθετους παιδαγωγικούς σχεδιασμούς.

Γι' αυτόν ακριβώς τον λόγο, η προδιαγραφή επιτρέπει προσαρμοσμένους \ENG{rollup rules}. Ο δημιουργός μπορεί να δηλώσει ποιο υποσύνολο παιδιών συνεισφέρει στο \ENG{rollup} μέσω κριτηρίων όπως \ENG{all}, \ENG{any}, \ENG{atLeastCount} ή \ENG{atLeastPercent}, να διατυπώσει λογικές συνθήκες πάνω στην κατάστασή τους και να καθορίσει ποια κατάσταση θα ανατεθεί στον γονέα. Με άλλα λόγια, το \ENG{rollup} μετατρέπει τη γονική δραστηριότητα σε υπολογιστικό αποτέλεσμα της συμπεριφοράς των παιδικών δραστηριοτήτων και όχι σε απλή διοικητική ομαδοποίηση.

Το ακόλουθο παράδειγμα παρουσιάζει έναν κανόνα \ENG{rollup} για απαίτηση καθολικής επιτυχίας.

{\selectlanguage{english}\fontencoding{T1}\selectfont

\begin{verbatim}
<imsss:rollupRule childActivitySet="all">
  <imsss:rollupConditions conditionCombination="any">
    <imsss:rollupCondition condition="satisfied"/>
  </imsss:rollupConditions>
  <imsss:rollupAction action="satisfied"/>
</imsss:rollupRule>
\end{verbatim}
}

Ο κανόνας αυτός χαρακτηρίζει τη γονική \ENG{activity} ως \ENG{satisfied} όταν ικανοποιείται η δηλωμένη συνθήκη για το σύνολο των παιδιών που συμμετέχουν στο \ENG{rollup}. Σε συνδυασμό με άλλους κανόνες, επιτρέπει τη μοντελοποίηση αυστηρών απαιτήσεων επιτυχίας, οριακών περιπτώσεων ή πολυεπίπεδων σεναρίων αξιολόγησης~\cite{ADL_SCORM2004_SN, IMS_SS_RollupModel}.

\subsubsection{Έλεγχοι \ENG{Navigation}}

Το \ENG{navigation} στο \ENG{SCORM} 2004 δεν είναι αυτόνομο σύστημα, αλλά το λειτουργικό πρόσωπο του \ENG{sequencing} απέναντι στον εκπαιδευόμενο και στο ίδιο το \ENG{SCO}. Τα αιτήματα πλοήγησης ενεργοποιούν διαδικασίες \ENG{sequencing}, οι οποίες με τη σειρά τους αποφασίζουν αν η αιτούμενη μετάβαση είναι επιτρεπτή και ποια δραστηριότητα θα παραδοθεί στη συνέχεια~\cite{ADL_SCORM2004_SN, SCORM_Navigation_Guide}.

Τα αιτήματα αυτά μπορούν να προέρχονται από τρεις πηγές. Πρώτον, από ενέργειες του εκπαιδευόμενου στα χειριστήρια πλοήγησης του \ENG{LMS}, όπως επιλογές στον \ENG{TOC} ή πλήκτρα \ENG{previous/next}. Δεύτερον, από το ίδιο το \ENG{SCO}, όταν αυτό θέτει μια τιμή στο \texttt{\ENG{adl.nav.request}} πριν τερματίσει. Τρίτον, από τις ίδιες τις ενέργειες που προβλέπουν οι \ENG{sequencing rules}, όταν για παράδειγμα ένας \ENG{post-condition rule} προκαλεί \ENG{continue} ή \ENG{exitAll}. Η ροή αυτών των αιτημάτων παρουσιάζεται στο Σχήμα~\ref{fig:navigation_controls}, ενώ ο Πίνακας~\ref{tab:navigation_events} συνοψίζει τα αντίστοιχα \ENG{navigation events} και τις επιπτώσεις τους~\cite{ADL_SCORM2004_SN, ADL_SCORM_Users_Guide}.

\begin{figure}[htbp]
\centering
% Figure 3.10: SCORM 2004 Navigation Controls and Request Flow
\begin{chapterfigurebox}
\begin{tikzpicture}[
    component/.style={rectangle, draw=blue!60, fill=blue!10, thick, minimum width=3.4cm, minimum height=1.1cm, align=center, rounded corners},
    source/.style={ellipse, draw=green!60, fill=green!10, thick, minimum width=2.9cm, minimum height=1.0cm, align=center},
    request/.style={rectangle, draw=orange!60, fill=orange!10, thick, minimum width=2.45cm, minimum height=0.85cm, align=center, rounded corners, font=\scriptsize},
    arrow/.style={->, >=stealth, thick}
]

% Navigation Request Sources
\node[source] (learner) at (-4.8,0) {\ENG{Learner}\\\ENG{Action}};
\node[source] (sco_nav) at (0,0) {\ENG{SCO} \ENG{Nav}\\\ENG{Request}};
\node[source] (seq_action) at (4.8,0) {\ENG{Sequencing}\\\ENG{Action}};

% Central Navigation Request Handler
\node[component, minimum width=9.2cm, minimum height=1.35cm] (nav_handler) at (0,-2.4) {\ENG{Navigation} \ENG{Request} \ENG{Handler}};

\draw[arrow] (learner) -- (nav_handler);
\draw[arrow] (sco_nav) -- (nav_handler);
\draw[arrow] (seq_action) -- (nav_handler);

\node[request] (continue_evt) at (-4.9,-5.2) {\ENG{continue}};
\node[request] (previous_evt) at (-1.6,-5.2) {\ENG{previous}};
\node[request] (choice_evt) at (1.6,-5.2) {\ENG{choice}};
\node[request] (exit_evt) at (4.9,-5.2) {\ENG{exit}};
\node[request] (exitAll_evt) at (-4.9,-6.8) {\ENG{exitAll}};
\node[request] (abandon_evt) at (-1.6,-6.8) {\ENG{abandon}};
\node[request] (suspendAll_evt) at (1.6,-6.8) {\ENG{suspendAll}};
\node[request] (abandonAll_evt) at (4.9,-6.8) {\ENG{abandonAll}};

\node[
    draw=gray!60,
    dashed,
    rounded corners,
    inner sep=0.45cm,
    fit=(continue_evt) (abandonAll_evt),
    label={[font=\scriptsize\bfseries, anchor=north west]north west:\ENG{Navigation} \ENG{requests}}
] (request_group) {};

\draw[arrow] (nav_handler) -- (request_group.north);

% Sequencing Process
\node[component, minimum width=8cm] (sequencing) at (0,-9.6) {\ENG{Sequencing} \ENG{Process}};
\draw[arrow] (request_group.south) -- (sequencing.north);

% Control Modes
\node[component, fill=purple!10, draw=purple!60, text width=11.4cm, minimum height=1.95cm] (control_modes) at (0,-12.6) {\ENG{Control} \ENG{modes} \ENG{and} \ENG{constraints}: \ENG{Choice}, \ENG{Flow}, \ENG{ForwardOnly}, \ENG{ChoiceExit}, \ENG{ConstrainChoice}, \ENG{PreventActivation}};

\draw[arrow, dashed, thick, red!60] (control_modes) -- node[right, font=\scriptsize] {\ENG{Restrict}} (sequencing);

% Annotations
\node[font=\scriptsize\itshape, text width=3.4cm, align=left, above=0.4cm of learner] {\ENG{TOC} \ENG{selection},\\\ENG{buttons}};
\node[font=\scriptsize\itshape, text width=3.4cm, align=center, above=0.4cm of sco_nav] {\ENG{adl.nav.request}};
\node[font=\scriptsize\itshape, text width=3.4cm, align=right, above=0.4cm of seq_action] {\ENG{Rule} \ENG{actions}};

\end{tikzpicture}
\end{chapterfigurebox}

\caption{Έλεγχοι \ENG{Navigation} και ροή \ENG{Request} στο \ENG{SCORM} 2004}
\label{fig:navigation_controls}
\end{figure}

\begin{table}[htbp]
\centering
\caption{\ENG{SCORM} 2004 \ENG{Navigation Events}}
\label{tab:navigation_events}
\begingroup
\footnotesize
\setlength{\tabcolsep}{3pt}
\renewcommand{\arraystretch}{1.16}
\begin{adjustbox}{center,max width=\textwidth,max totalheight=0.86\textheight}
\begin{tabular}{|p{3cm}|p{3cm}|p{7.5cm}|}
\hline
\ENG{Navigation} \ENG{Event} & \ENG{Source} & \ENG{Description} \ENG{and} \ENG{Sequencing} \ENG{Effect} \\
\hline
\hline
\texttt{\ENG{continue}} & \ENG{Learner}, \ENG{SCO}, \ENG{Sequencing} & \ENG{Requests} \ENG{delivery} \ENG{of} \ENG{the} \ENG{logically} "\ENG{next}" \ENG{activity} \ENG{in} \ENG{a} \ENG{flow-enabled} \ENG{cluster}. \ENG{Triggers} \texttt{\ENG{Flow}} \ENG{sequencing} \ENG{subprocess}. \ENG{If} \ENG{no} \ENG{next} \ENG{activity} \ENG{exists}, \ENG{may} \ENG{exit} \ENG{parent} \ENG{or} \ENG{course}. \\
\hline
\texttt{\ENG{previous}} & \ENG{Learner}, \ENG{SCO}, \ENG{Sequencing} & \ENG{Requests} \ENG{delivery} \ENG{of} \ENG{the} \ENG{logically} "\ENG{previous}" \ENG{activity} \ENG{in} \ENG{a} \ENG{flow-enabled} \ENG{cluster}. \ENG{Triggers} \ENG{reverse} \texttt{\ENG{Flow}} \ENG{sequencing} \ENG{subprocess}. \ENG{Restricted} \ENG{by} \texttt{\ENG{forwardOnly}} \ENG{control} \ENG{mode}. \\
\hline
\texttt{\ENG{choice}} & \ENG{Learner}, \ENG{SCO} & \ENG{Requests} \ENG{delivery} \ENG{of} \ENG{a} \ENG{specific} \ENG{activity} \ENG{identified} \ENG{by} \ENG{its} \texttt{\ENG{identifier}}. \ENG{Valid} \ENG{only} \ENG{if} \texttt{\ENG{choice}} \ENG{control} \ENG{mode} \ENG{is} \ENG{enabled} \ENG{for} \ENG{the} \ENG{parent} \ENG{cluster}. \ENG{Subject} \ENG{to} \texttt{\ENG{constrainChoice}} \ENG{and} \texttt{\ENG{preventActivation}} \ENG{restrictions}. \\
\hline
\texttt{\ENG{exit}} & \ENG{Learner}, \ENG{SCO}, \ENG{Sequencing} & \ENG{Terminates} \ENG{the} \ENG{current} \ENG{activity} \ENG{and} \ENG{bubbles} \ENG{exit} \ENG{request} \ENG{to} \ENG{parent}. \ENG{Triggers} \ENG{parent}'\ENG{s} \ENG{exit} \ENG{sequencing} \ENG{rules}. \ENG{May} \ENG{result} \ENG{in} \ENG{delivery} \ENG{of} \ENG{next} \ENG{activity} \ENG{or} \ENG{course} \ENG{termination}. \\
\hline
\texttt{\ENG{exitAll}} & \ENG{Learner}, \ENG{SCO}, \ENG{Sequencing} & \ENG{Terminates} \ENG{the} \ENG{current} \ENG{activity} \ENG{and} \ENG{all} \ENG{ancestors} \ENG{up} \ENG{to} \ENG{the} \ENG{root}. \ENG{Causes} \ENG{course} \ENG{to} \ENG{exit}. \ENG{Learner} \ENG{progress} \ENG{and} \ENG{status} \ENG{are} \ENG{persisted}. \\
\hline
\texttt{\ENG{abandon}} & \ENG{Learner}, \ENG{SCO} & \ENG{Terminates} \ENG{the} \ENG{current} \ENG{activity} \ENG{without} \ENG{recording} \ENG{its} \ENG{status} \ENG{or} \ENG{triggering} \ENG{rollup}. \ENG{Similar} \ENG{to} \texttt{\ENG{exit}} \ENG{but} \ENG{discards} \ENG{current} \ENG{activity}'\ENG{s} \ENG{tracking} \ENG{data}. \ENG{Rarely} \ENG{used} \ENG{in} \ENG{practice}. \\
\hline
\texttt{\ENG{abandonAll}} & \ENG{Learner}, \ENG{SCO} & \ENG{Terminates} \ENG{the} \ENG{current} \ENG{activity} \ENG{and} \ENG{all} \ENG{ancestors} \ENG{without} \ENG{recording} \ENG{status}. \ENG{Causes} \ENG{course} \ENG{to} \ENG{exit} \ENG{without} \ENG{rollup}. \ENG{Extreme} \ENG{edge} \ENG{case}; \ENG{rarely} \ENG{supported} \ENG{fully}. \\
\hline
\texttt{\ENG{suspendAll}} & \ENG{Learner}, \ENG{SCO} & \ENG{Suspends} \ENG{the} \ENG{course}, \ENG{preserving} \ENG{all} \ENG{tracking} \ENG{data} \ENG{and} \ENG{current} \ENG{activity} \ENG{state}. \ENG{Next} \ENG{launch} \ENG{resumes} \ENG{from} \ENG{suspended} \ENG{state}. \ENG{Sets} \texttt{\ENG{cmi.exit}} \ENG{to} \texttt{"\ENG{suspend}"}. \\
\hline
\texttt{\_\ENG{none}\_} & \ENG{SCO} (\ENG{default}) & \ENG{Indicates} \ENG{no} \ENG{explicit} \ENG{navigation} \ENG{request}. \ENG{Sequencer} \ENG{applies} \ENG{post-condition} \ENG{rules} \ENG{to} \ENG{determine} \ENG{subsequent} \ENG{action} (\ENG{may} \ENG{result} \ENG{in} \ENG{continue}, \ENG{exit}, \ENG{or} \ENG{retry} \ENG{based} \ENG{on} \ENG{rules}). \\
\hline
\end{tabular}
\end{adjustbox}
\endgroup

\end{table}

Οι \ENG{sequencing control modes} καθορίζουν το εύρος των επιτρεπτών μεταβάσεων μέσα στα παιδιά μιας \ENG{cluster activity}. Το \ENG{choice} επιτρέπει πλοήγηση βασισμένη στον \ENG{TOC}, το \ENG{flow} ενεργοποιεί ευθύγραμμη πλοήγηση τύπου \ENG{previous/next}, το \ENG{forward only} απαγορεύει την επιστροφή προς τα πίσω και το \ENG{choice exit} ελέγχει αν ο εκπαιδευόμενος μπορεί να εγκαταλείψει μια δραστηριότητα μέσω επιλογής από τον πίνακα περιεχομένων. Πάνω σε αυτούς τους βασικούς τρόπους λειτουργίας προστίθενται έλεγχοι όπως το \ENG{constrain choice}, που περιορίζει τις επιτρεπτές επιλογές στις λογικά διαθέσιμες επόμενες δραστηριότητες, και το \ENG{prevent activation}, που αποτρέπει το \ENG{deep linking} προς παιδιά χωρίς προηγούμενη είσοδο στον γονέα. Έτσι το \ENG{SCORM} 2004 μπορεί να υποστηρίξει τόσο αμιγώς γραμμικά όσο και περισσότερο ανοικτά ή υβριδικά μοντέλα πλοήγησης~\cite{ADL_SCORM2004_SN, IMS_SS_ControlModes}.

Παράλληλα, το στοιχείο \texttt{<\ENG{imsss:hideLMSUI}>} επιτρέπει στο περιεχόμενο να δηλώσει ότι συγκεκριμένα χειριστήρια του \ENG{LMS}, όπως \ENG{previous}, \ENG{continue}, \ENG{exit}, \ENG{exitAll}, \ENG{abandon}, \ENG{abandonAll} ή \ENG{suspendAll}, πρέπει να αποκρύπτονται. Η δυνατότητα αυτή δεν καταργεί τον έλεγχο του \ENG{LMS}, αλλά του επιτρέπει να συντονίσει τη διεπαφή με την εσωτερική λογική του μαθήματος, αποφεύγοντας αντιφατικά σήματα προς τον εκπαιδευόμενο όταν το ίδιο το \ENG{SCO} παρέχει εξειδικευμένη πλοήγηση~\cite{ADL_SCORM2004_SN, SCORM_Navigation_Guide}.

\subsubsection{Καθολικά \ENG{Objectives} και \ENG{Objective Mapping}}

Ένα από τα ισχυρότερα αλλά συχνά λιγότερο αξιοποιημένα χαρακτηριστικά του \ENG{SCORM} 2004 είναι η υποστήριξη καθολικών \ENG{objectives}. Πρόκειται για στόχους οι οποίοι δεν περιορίζονται στη διάρκεια ζωής μιας μόνο δραστηριότητας ή ενός μόνο πακέτου, αλλά μπορούν να χρησιμοποιηθούν ως κοινά σημεία αναφοράς ανάμεσα σε πολλαπλές \ENG{activities} ή ακόμη και σε διαφορετικά μαθήματα. Με τον τρόπο αυτό καθίσταται εφικτός ο έλεγχος προαπαιτουμένων με βάση εξωτερικές επιδόσεις, η διατήρηση κοινής κατάστασης δεξιοτήτων ή \ENG{competencies} και η υποστήριξη \ENG{test-out} σεναρίων, όπου επιτυχία σε μια αρχική αξιολόγηση επηρεάζει τη διαθεσιμότητα επόμενων δραστηριοτήτων.

Το \ENG{objective mapping} είναι ο μηχανισμός μέσω του οποίου μια τοπική \ENG{activity} συσχετίζει τους δικούς της στόχους με ένα καθολικό \ENG{objective}. Η χαρτογράφηση αυτή δηλώνει το τοπικό και το καθολικό αναγνωριστικό, καθώς και τα δικαιώματα ανάγνωσης και εγγραφής για στοιχεία όπως το \ENG{satisfaction status}, το \ENG{measure}, οι βαθμολογίες \ENG{raw/min/max}, το \ENG{completion status} και το \ENG{progress measure}. Ιδίως στην \ENG{4th Edition}, η ενίσχυση αυτού του μηχανισμού καθιστά το \ENG{SCORM} 2004 ικανό να μοντελοποιήσει σενάρια μάθησης που εκτείνονται πέρα από τα όρια μιας μεμονωμένης εκτέλεσης~\cite{ADL_SCORM2004_4thEd, ADL_SCORM_Users_Guide}.

\subsubsection{Στρατηγικές \ENG{Sequencing} και περιπτώσεις χρήσης}

Η εκφραστικότητα του \ENG{SCORM} 2004 γίνεται ιδιαίτερα εμφανής όταν εξεταστούν οι στρατηγικές που μπορεί να υποστηρίξει. Με κατάλληλους \ENG{precondition rules} μπορούν να υλοποιηθούν αυστηρά προαπαιτούμενα και αλληλεξαρτήσεις ανάμεσα σε ενότητες. Με συνδυασμό \ENG{flow}, απενεργοποιημένου \ENG{choice} και ενεργειών όπως \ENG{stopForwardTraversal} μπορεί να επιβληθεί αυστηρά γραμμική πορεία. Οι \ENG{post-condition rules} επιτρέπουν κύκλους επανεκπαίδευσης και \ENG{remediation}, ενώ τα καθολικά \ENG{objectives} καθιστούν εφικτά \ENG{test-out} σενάρια στα οποία η επιτυχία σε προαξιολόγηση επιτρέπει την παράκαμψη τμημάτων του μαθήματος. Επιπλέον, μέσω \ENG{selection controls} μπορούν να δημιουργηθούν τυχαίες επιλογές από δεξαμενές ερωτήσεων, ενώ ο συνδυασμός \ENG{rollup rules} και κατάστασης στόχων υποστηρίζει προσαρμοστικές μαθησιακές διαδρομές.

Η ισχύς αυτή, ωστόσο, συνοδεύεται από σημαντική πολυπλοκότητα. Η προδιαγραφή \ENG{IMS Simple Sequencing} περιλαμβάνει εκτενές \ENG{pseudocode} που ορίζει με ακρίβεια τους αλγορίθμους επιλογής δραστηριότητας, επεξεργασίας αιτημάτων πλοήγησης και \ENG{rollup}. Συνεπώς, η σωστή υλοποίηση ενός πλήρους \ENG{SCORM} 2004 \ENG{LMS} ή \ENG{sequencing engine} απαιτεί πολύ υψηλότερο επίπεδο τεχνικής ωριμότητας σε σχέση με το \ENG{SCORM} 1.2, γεγονός που εξηγεί και την άνιση υιοθέτηση του προτύπου στην πράξη~\cite{IMS_SS_BehaviorModel, ADL_SCORM2004_SN}.

\subsection{Σύνοψη}
\label{subsec:scorm2004_summary}

Το \ENG{SCORM} 2004 αποτελεί ώριμη αλλά απαιτητική εξέλιξη του \ENG{SCORM} 1.2. Διατηρεί τη βασική αρχή του διαλειτουργικού πακεταρίσματος και της \ENG{browser-based} επικοινωνίας με το \ENG{LMS}, αλλά την επεκτείνει με αυστηρότερο \ENG{runtime}, αναλυτικότερο μοντέλο δεδομένων και, κυρίως, με τυποποιημένο μηχανισμό \ENG{sequencing and navigation}. Ως αποτέλεσμα, το πρότυπο μπορεί να υποστηρίξει σύνθετες παιδαγωγικές στρατηγικές, λεπτομερέστερη παρακολούθηση και περισσότερο προσαρμοστικές μαθησιακές εμπειρίες.

Η ίδια ακριβώς ισχύς αποτελεί και το βασικό του μειονέκτημα. Η υλοποίηση του \ENG{SCORM} 2004, ιδίως στο επίπεδο του \ENG{sequencing}, είναι σημαντικά πιο πολύπλοκη από εκείνη του \ENG{SCORM} 1.2, τόσο για τους δημιουργούς περιεχομένου όσο και για τους κατασκευαστές \ENG{LMS}. Η ιστορική συνύπαρξη τεχνικής υπεροχής και περιορισμένης επιχειρησιακής υιοθέτησης αποτυπώνει εύγλωττα το κεντρικό δίλημμα του προτύπου: όσο αυξάνεται η εκφραστικότητα και η παιδαγωγική δύναμη, τόσο αυξάνεται και το κόστος πλήρους και ορθής υλοποίησης.
